\documentclass[11pt,a4paper]{article}
\usepackage[T1]{fontenc}
\usepackage[utf8]{inputenc}
\usepackage{lmodern}
\usepackage{amsmath,amssymb}
\usepackage[margin=25mm]{geometry}
\usepackage[hidelinks]{hyperref}
\hypersetup{pdftitle={Shared radiation and closure under recording},pdfauthor={navstokgap research working notes}}
\newcommand{\tightlist}{\setlength{\itemsep}{0pt}\setlength{\parskip}{0pt}}
\setlength{\emergencystretch}{3em}
\setcounter{secnumdepth}{0}
\begin{document}
\hypertarget{shared-radiation-and-closure-under-recording}{%
\section{Shared radiation and closure under
recording}\label{shared-radiation-and-closure-under-recording}}

\textbf{Physical adaptive-gain result, 2026-09-30 (GPT-6.1 Sol and GPT-6
Astra; refereed).} Section 8.10 realizes feedback by a Hamiltonian
coupling directly to stored physical memory coordinates, retaining the
controller recoil on their conjugates. Its conditional body posterior is
exactly the Gaussian posterior of the gain sequence frozen at the
observed record. Closure and the complete-record continuity bound (77)
hold uniformly over fixed partitions and bounded smooth gains. The same
continuity holds at zero action. Adaptive timing and refinement
compatibility between different gain rules remain subsequent tests.

\textbf{Associative terminal-channel result, 2026-09-30 (GPT-6.1 Sol and
GPT-6 Astra; written, internally checked).} Sections 8.7--8.9 retain
every physical copy kick in an associative Gaussian body--record
descriptor. Three equal cells have the rank-two relative residue (69).
For any fixed duration and precision density, the projected terminal law
converges at the explicit mesh rate (73) as the mesh tends to zero,
independently of the partition schedule. The retained record is one
weighted terminal projection; a mechanical trajectory and the limit of
the entire growing record are separate constructions. The posterior
floor survives this projection and its limit when the joint preparation
scale is positive. The same channel convergence holds on the zero-action
branch.

\textbf{Two-time blocking result, 2026-09-29 (GPT-6.1 Sol and GPT-6
Astra; written, internally checked).} Sections 8.4--8.6 construct two
physical memory copies at distinct times and compare their projected
joint body--record law to one effective copy at the same terminal cut.
The relative kick survives as an explicit rank-one noise, with the
split-uniform statistical defect (62). A fixed precision density gives
the complete terminal record a partition-uniform force-information bound
(63), even on the zero-action branch. Sections 8.7--8.9 develop
associative three-cell blocking; the positive preparation scale remains
an input.

\textbf{Hamiltonian memory result, 2026-09-29 (GPT-6.1 Sol and GPT-6
Astra; written, internally checked).} Section 8 constructs the kick
monitor as a coupling to a physical memory. It changes the first stored
imprecision, leaving the conditional noise determinant (52) at least
\(\kappa^2\). For a general affine Hamiltonian apparatus, the complete
terminal memory coordinates imply the Gaussian posterior floor from the
initial joint covariance restriction (22). Linear blocking of those
records preserves the floor. This derives the readout cost in the stated
class; its initial positive phase-covariance scale remains a premise.

\textbf{Adaptive result, 2026-09-29 (GPT-6.1 Sol; written, unrefereed).}
Section 7 gives the exact Hellinger affinity of a finite adaptive linear
Gaussian record in terms of its branchwise information. It includes
retained memory and indirect noise observations. A uniform branchwise
budget controls both force distinguishability and, with the stated
matching covariance assumption, refinement of the complete record.
Monitoring a saturated pointer kick requires fresh hidden variance
\(b_0^2/(b_0+r^2)\) in the explicit repair (§7.2). This quantifies a
possible innovation mechanism; the physical law and its positive scale
remain to be derived.

\textbf{Further result, 2026-09-29 (GPT-6.1 Sol; written, unrefereed).}
In the two-pointer model, allow correlated Gaussian readout error and
kick \((N,D)\) and side records \(S\) of that noise. The body's
conditional area remains at least \(\kappa\) exactly when
\(\det\operatorname{Cov}(N,D\mid S)\ge\kappa^2\) (§6.2), under the
independence assumptions stated there. Even a noisy side observation of
the kick breaks a marginally saturated recording law. Uniform continuity
of the complete Gaussian record is equivalent to bounded full-record
Fisher information (§6.1); continuity of each apparatus separately does
not give that bound. The finite Lean results suggest testing both
common-mode cancellation and refinement error in the record's noise norm
(§6.3).

\textbf{Result, 2026-09-27 (GPT-6 Astra).} A common Gaussian radiation
background permits arbitrarily small two-pointer posteriors in a linear
classical model with the pulse couplings and retained coordinate records
of Theorem I. This holds for \textbf{every finite joint covariance},
including all body--pointer and pointer--pointer correlations. At fixed
ultraviolet cutoff, radiation damping and continuing exposure to the
background preserve the counterexample for sufficiently short finite
pulses (Theorems 1 and 2). At fixed cutoff the infimum of the body's
posterior area is zero. For ideal resonant marginal covariances of area
\(\kappa\), finite gains \(2\) and \(4\) already give
\(\sqrt{\det\Sigma_{\rm post}}\le3\kappa/4\).

\textbf{After refereeing (Fable, two passes, 2026-09-27; all theorems
accepted).} The shared background's own prior already satisfies the
Gaussian restriction (22): each bath mode sits on its boundary and the
flow is symplectic. So Theorems 1--2 break closure through the readout
class alone, a passive, noise-free retained reading. Any Hamiltonian
readout of a pointer coordinate kicks its conjugate; for a resonant
device in the background, imprecision times kick equals the device area
\(\kappa\), which is Proposition 3's law, and nesting pointers only
relocates the passive readout. At \(\kappa=\hbar/2\) the recording law
(20)--(22) is Gaussian quantum measurement theory term by term (the
Simon--Mukunda--Dutta state condition, general-dyne measurements, the
Giedke--Cirac update); Bartlett, Rudolph and Spekkens prove the
operational equivalence. The minimal premise that restores closure is
therefore a bound on every terminal classical readout, with retained
pairs read only through couplings. It is a restriction on observation,
and the classical dynamics of the background cannot supply it.

The decisive instrument premise is classical storage of a pointer
coordinate without an obligatory conjugate disturbance when that pointer
is subsequently reused. An explicit Gaussian recording law restores
closure at \(\zeta=\kappa\): readout imprecision and the added conjugate
disturbance must satisfy a product bound \(\kappa^2\), with a
corresponding restriction on every retained memory (Proposition 3).
Deriving that law from electrodynamics remains the physical task. The
stationary spectrum alone leaves it unspecified.

\textbf{Scope.} The cutoff model below supplies the exact
finite-duration statement. Its background has mode energy
\(\kappa\omega\) below the cutoff and a specified regulator above it. A
laboratory frequency cutoff adds a preferred frame. The resonant
approximation gives the simpler constant \(3\kappa/4\); the proof holds
for arbitrary finite covariance, with or without those simplifications.
The unregulated full momentum variance is divergent, as stressed in the
revised \href{sed-zeta-radiation-link.md}{SED link}, §2 (local
full-read). The cutoff stays fixed throughout. The controls and
classical readouts are explicitly granted model operations; an
electrodynamic construction of a complete measuring device would have to
justify or exclude them.

\hypertarget{the-common-field-its-damping-and-its-correlations}{%
\subsection{1. The common field, its damping and its
correlations}\label{the-common-field-its-damping-and-its-correlations}}

Use SI units, \(\kappa>0\), speed \(c\), and three harmonically bound
dipole coordinates \(q_i,p_i\), \(i=0,1,2\), with masses \(m_i>0\),
frequencies \(\omega_i>0\), charges \(e_i\ne0\), fixed centres
\(\mathbf r_i\), and unit directions \(\mathbf n_i\). The body has index
\(0\). The traps and external quadratic controls define the linear
approximation. The incident Gaussian field has the zero-point spectrum

\[u_0(\omega)=\frac{\kappa\omega^3}{\pi^2c^3},\qquad
\mathcal E_0(\omega)=\kappa\omega. \tag{1}\]

Introduce a real common form factor \(f_\Lambda(\omega)\), equal to one
on the frequencies of interest and zero for \(\omega>\Lambda\), with
\(\Lambda<\infty\). This can regulate the incident force itself or the
dipole coupling to the field. In the latter convention the full incident
field retains (1), while the apparatus response selects a frame and a
frequency band. All covariances below refer to the regulated response.

The angular correlation matrix is

\[K_{ij}(\omega)=\frac{3}{8\pi}\int_{S^2}
\left[\mathbf n_i\!\cdot\!\mathbf n_j
 -(\mathbf n_i\!\cdot\!\widehat{\mathbf k})
  (\mathbf n_j\!\cdot\!\widehat{\mathbf k})\right]
\cos\!\left(\frac{\omega}{c}\widehat{\mathbf k}\!\cdot\!
 (\mathbf r_i-\mathbf r_j)\right)d\widehat{\mathbf k}. \tag{2}\]

It is positive semidefinite and \(K_{ii}=1\): the transverse
polarization sum is a Gram matrix, and the integral for a unit direction
is \(8\pi/3\). Thus the same field gives both diagonal and cross force
covariances. Define its matrix spectral density by

\[J_{ij}(\omega)=\frac{e_i e_j\omega^3}{6\pi\epsilon_0c^3}
 f_\Lambda(\omega)^2K_{ij}(\omega),\]
\[C_F(t-s)=\mathbb E[F(t)F(s)^{\sf T}]
 =\frac{2\kappa}{\pi}\int_0^\Lambda J(\omega)
       \cos[\omega(t-s)]\,d\omega. \tag{3}\]

For example
\(C_{F,ii}(0)=e_i^2\kappa\int_0^\Lambda \omega^3f_\Lambda^2d\omega/(3\pi^2\epsilon_0c^3)\).
The normalization follows directly from
\(u_0=\epsilon_0\langle |\mathbf E|^2\rangle\) per frequency, isotropy,
and \(F_i=e_i E_i\).

Keep radiation damping through a passive oscillator bath with the same
\(J\), including the cross damping. Its causal friction kernel is

\[\mu(t)=\frac2\pi\int_0^\Lambda
       \frac{J(\omega)}{\omega}\cos(\omega t)\,d\omega,
\qquad t\ge0. \tag{4}\]

An explicit realization uses harmonic bath coordinates coupled linearly
to the \(q_i\), with the quadratic counterterm obtained by completing
each bath potential square. Choose the bath coupling matrix so that
\(J(\omega)=\pi L(\omega)L(\omega)^{\sf T}/(2\omega)\); its existence
follows from positivity of (2). Give a free bath coordinate of frequency
\(\omega\) the variances \(\kappa/\omega\) and \(\kappa\omega\) for
coordinate and momentum. Its energy is \(\kappa\omega\). Eliminating it
gives (3) and (4). The counterterm keeps the prescribed trap stiffness
positive; the finite-cutoff bath Hamiltonian has positive quadratic
energy.

In this realization \(\dot q=\partial H_{\rm sys}(t)/\partial p\), and

\[\dot p=-\frac{\partial H_{\rm sys}(t)}{\partial q}
 -\int_0^t\mu(t-s)\dot q(s)\,ds+F(t). \tag{5}\]

An initial-slip term, or the stationary prehistory, is included in
\(F(t)\) with its correlations with the initial system. Independent
system--bath initial data are unnecessary below. Before the controls,
the diagonal absorptive part is the usual radiation term
\(m_i\tau_i\omega^3 f_\Lambda^2\), where

\[\tau_i=\frac{e_i^2}{6\pi\epsilon_0m_i c^3}.
\tag{6}\]

This finite-cutoff memory equation retains the dispersive part as well.
It avoids using a third-order local radiation equation during
arbitrarily fast controls. Static dipole interactions, when retained,
are additional finite entries of the stiffness matrix; assume stability
of the chosen trapped system.

\textbf{Tracking the prior.} Set

\[Q_i=\sqrt{m_i\omega_i}\,q_i,\qquad
P_i=\frac{p_i}{\sqrt{m_i\omega_i}},\qquad
Z=(Q,P,Y,\Pi,Z_2,\Pi_2)^{\sf T}. \tag{7}\]

Here \((Q,P)=(Q_0,P_0)\), \((Y,\Pi)=(Q_1,P_1)\) and
\((Z_2,\Pi_2)=(Q_2,P_2)\). Every coordinate has units
\(\sqrt{\rm action}\), every pair remains canonical, and the body's
covariance determinant equals that in \((q_0,p_0)\). For a stable
stationary response, let \(H(\omega)\) be the six-by-three transfer
matrix from physical forces to these variables. The complete prior
covariance is

\[V=\operatorname{Re}\int_0^\Lambda
 H(\omega)\frac{2\kappa J(\omega)}{\pi}
 H(\omega)^\dagger\,d\omega. \tag{8}\]

The susceptibility in \(H\) includes (4); momentum rows before the
controls are \(-i m_i\omega\) times the coordinate rows before scaling.
Any undamped initially populated normal modes add their own covariance.
Formula (8) displays the cross correlations explicitly, and the theorems
use the full \(V\) of (8), cross correlations included. Its dispersive
cross part and the constant \(K\) of (14) are cutoff constants
(\(\sigma_P^2\) grows with \(\Lambda\)). The radiation time here is the
SI \(\tau=e^2/(6\pi\epsilon_0mc^3)\); the SED link note uses the
Gaussian-unit form.

For an isolated weakly damped resonance the ideal narrow-resonance limit
gives \(\operatorname{Var}Q_i=\operatorname{Var}P_i=\kappa\). A finite
spectral window and finite damping introduce corrections; the exact
finite-cutoff statements below allow all of them. In particular, a
literal window of only a fixed number of linewidths captures only a
fraction of the resonant Lorentzian. Access to a temporally filtered
variable also requires a recording prescription; the instantaneous pulse
proof is applied to canonical cutoff variables, or to canonical
oscillator amplitudes in the stated resonant model.

\hypertarget{theorem-1-two-correlated-pointers-permit-arbitrarily-small-area}{%
\subsection{2. Theorem 1: two correlated pointers permit arbitrarily
small
area}\label{theorem-1-two-correlated-pointers-permit-arbitrarily-small-area}}

Assume \(Z\) has a centered Gaussian law with finite covariance \(V\);
known means can be subtracted. Grant two quadratic Hamiltonian pulses,
their integrated Hamiltonians being

\[\mathcal H_1=GQ\Pi,\qquad
\mathcal H_2=B\Pi\Pi_2,\qquad G,B>0. \tag{9}\]

Both have action units. Their Hamilton equations are linear. The first
pulse sends \(Y\mapsto Y+GQ\) and \(P\mapsto P-G\Pi\). Read and store
\(R_1=Y+GQ\) before the second pulse. That pulse sends
\(Z_2\mapsto Z_2+B\Pi\) and \(Y\mapsto Y+B\Pi_2\); it leaves the body
and \(\Pi\) fixed. Read \(R_2=Z_2+B\Pi\). The final body is

\[X=(Q,P-G\Pi)^{\sf T},\qquad R=(R_1,R_2)^{\sf T}. \tag{10}\]

The readouts are classical passive coordinate observations. The first
stored number survives the later change of \(Y\). Small independent
Gaussian readout errors can also be allowed, as in §3.

\textbf{Theorem 1.} Under these instrument assumptions, for every such
\(V\),

\[\boxed{\sqrt{\det\operatorname{Cov}(X\mid R)}
\le\frac{\sigma_Y\sigma_P}{G}
       +\frac{\sigma_Y\sigma_{Z_2}}{B}.} \tag{11}\]

Here \(\sigma_U^2=\operatorname{Var}U\) in the initial joint state. In
particular the posterior area has infimum zero as \(G,B\to\infty\). All
cross correlations in \(V\) are allowed, including perfect ones.

\emph{Proof.} Estimate the final body by

\[\widehat X=(R_1/G,-GR_2/B)^{\sf T}.\]

The errors are the exact random-variable identities

\[X-\widehat X=(-Y/G,\ P+GZ_2/B)^{\sf T}. \tag{12}\]

Thus the first pointer's recoil \(\Pi\) cancels, including the part
correlated with the body or either pointer. The conditional Gaussian
covariance is independent of the outcome. Conditional expectation
minimizes the error covariance in the positive-semidefinite order, so it
is bounded above by the covariance of (12). In particular,

\[\operatorname{Var}(Q\mid R)\le\sigma_Y^2/G^2,\qquad
\operatorname{Var}(P-G\Pi\mid R)
\le(\sigma_P+G\sigma_{Z_2}/B)^2.\]

The second bound is Cauchy--Schwarz, with either sign of
\(\operatorname{Cov}(P,Z_2)\). Bounding a two-by-two determinant by the
product of its diagonal entries proves (11). \(\square\)

For direct use of a specified common-field response, introduce

\[A=\begin{pmatrix}1&0&0&0&0&0\\0&1&0&-G&0&0\end{pmatrix},
\qquad
L=\begin{pmatrix}G&0&1&0&0&0\\0&0&0&B&1&0\end{pmatrix}.\]

The exact answer, with every entry from (8) retained, is

\[\Sigma_{\rm post}=AVA^{\sf T}
 -AVL^{\sf T}(LVL^{\sf T})^{-1}LVA^{\sf T}. \tag{13}\]

Use a pseudoinverse on the supported record subspace when necessary. For
example the two record covariances include

\[\operatorname{Cov}(R_1,R_2)
 =G V_{QZ_2}+GB V_{Q\Pi}+V_{YZ_2}+B V_{Y\Pi}.\]

The covariance of the error vector in (12) has off-diagonal entry
\(-V_{YP}/G-V_{YZ_2}/B\) and second diagonal entry
\(V_{PP}+2G V_{PZ_2}/B+G^2V_{Z_2Z_2}/B^2\). These formulas keep every
cross entry of \(V\). Equation (11) bounds their most unfavourable
possible effect.

\textbf{Explicit constants.} If
\(\sigma_Y^2,\sigma_P^2, \sigma_{Z_2}^2\le K\kappa\), put \(B=G^2\).
Then

\[\sqrt{\det\Sigma_{\rm post}}
 \le K\kappa(G^{-1}+G^{-2}). \tag{14}\]

Every finite-cutoff prior has a finite such \(K\). Taking \(G>2K+1\)
gives a strict bound below \(\kappa\). For ideal resonant marginals,
\(K=1\), and \(G=2\), \(B=4\) give \(3\kappa/4\), regardless of the
cross correlations. As an algebraic example only, for \(V=\kappa I_6\),

\[\Sigma_{\rm post}=
\begin{pmatrix}
\kappa/(1+G^2)&0\\
0&\kappa[1+G^2/(1+B^2)]
\end{pmatrix}. \tag{15}\]

With \(G=2\), \(B=4\), its area is \(\kappa\sqrt{21/85}\). Independence
is used only in this example. The body momentum's unconditional spread
grows with \(G\); the completed record estimates its kick. The theorem
concerns posterior concentration and leaves the unconditional
disturbance cost to its own accounting.

\hypertarget{theorem-2-a-finite-duration-version-with-radiation-damping}{%
\subsection{3. Theorem 2: a finite-duration version with radiation
damping}\label{theorem-2-a-finite-duration-version-with-radiation-damping}}

Fix \(\Lambda\), the initial finite Gaussian covariances, and finite
gains \(G,B\). Apply the two pulses over successive intervals of total
length \(T\), using smooth fixed pulse shapes rescaled by \(1/T\). Keep
(5) active throughout and retain the first reading at the interval
boundary. Let \(X_T,R_T\) denote the resulting final body and records.
Write their deviations from the ideal expressions (10) as
\(\Delta Q,\Delta P, \Delta R_1,\Delta R_2\), and set

\[d_Q(T)=\|\Delta Q-\Delta R_1/G\|_2,\qquad
d_P(T)=\|\Delta P+G\Delta R_2/B\|_2. \tag{16}\]

Norms are centered root-mean-square norms. They include readout errors
when present. The initial field and body can be correlated.

\textbf{Theorem 2.} In the finite-cutoff passive linear model (3)--(5),
with the classical readouts of §2, \(d_Q(T),d_P(T)\to0\) as
\(T\downarrow0\). Moreover

\[\boxed{\sqrt{\det\operatorname{Cov}(X_T\mid R_T)}}
\ \le\left(\frac{\sigma_Y}{G}+d_Q(T)\right)
\left(\sigma_P+\frac{G\sigma_{Z_2}}B+d_P(T)\right). \tag{17}\]

Consequently every strict violation in (11) persists for sufficiently
small positive \(T\). At \(K=1\), \(G=2\), \(B=4\), the explicit
sufficient error budget \(d_Q,d_P\le\sqrt\kappa/10\) gives area at most
\(24\kappa/25<\kappa\).

\emph{Proof.} At fixed cutoff \(C_F(0)\), \(\mu\) and \(\mu'\) are
finite. In particular

\[\left\|\int_0^T F_i(t)\,dt\right\|_2
 \le T\sqrt{C_{F,ii}(0)}.\]

Rapid control-induced changes of \(q\) leave the memory term bounded:
integration by parts gives

\[\int_0^t\mu(t-s)\dot q(s)\,ds
 =\mu(0)q(t)-\mu(t)q(0)
       +\int_0^t\mu'(t-s)q(s)\,ds. \tag{18}\]

For the control-only equation each pulse is a finite shear. In the
scaled coordinates its propagator and inverse have norms bounded,
uniformly in \(T\), by \((1+G)(1+B)\). Variation of constants, (18), and
the finite force second moment give a uniform mean-square bound on
\(q,p\) for sufficiently small \(T\) by the integral Gronwall
inequality. Integrating the remaining trap, bath and memory drifts then
makes the endpoint and first-read errors \(O(T)\), with constants
depending on \(\Lambda,G,B\), the initial covariances and trap
parameters. This proves the asserted convergence, while keeping
radiation damping in the equation. The same conclusion holds for any
supplied resonant Markov model with finite diffusion, with stochastic
increments \(O(\sqrt T)\).

Apply the estimator of §2 to \(R_T\). Its two errors equal (12) plus the
residuals in (16). The triangle inequality in \(L^2\) gives (17), and
Gaussian conditioning again makes the posterior covariance
outcome-independent. Finally \((1/2+1/10)(3/2+1/10)=24/25\). \(\square\)

\textbf{Correlations during the controls.} More explicitly, the
four-vector \(O_T=(X_T,R_T)\) has the form \(M_T Z+\eta_T\). If
\(W_T=\operatorname{Cov}(Z,\eta_T)\) and
\(N_T=\operatorname{Cov}(\eta_T)\), then

\[\operatorname{Cov}(O_T)=M_TVM_T^{\sf T}
 +M_TW_T+W_T^{\sf T}M_T^{\sf T}+N_T. \tag{19}\]

For a force-response kernel \(D_T(t)\) the driven part is

\[N_T=\int_0^T\!\int_0^T
D_T(t)C_F(t-s)D_T(s)^{\sf T}\,dt\,ds,\]

with any prehistory terms included in \(\eta_T\). Thus both the common
increments and their correlation with the initial state are retained.
Equation (17) uses an \(L^2\) bound, so (17) holds with every cross term
of (19) retained. Explicitly, with
\(\Lambda_*=\sqrt{\pi\omega_0/\tau}\), the budget
\(d_P\le\sqrt\kappa/10\) needs \(T\Lambda\le0.14\,(\Lambda_*/\Lambda)\)
from the force noise and \(T\Lambda\le0.15\,(\Lambda_*/\Lambda)^2\) from
memory: the pulses must outrun every retained field mode, and without a
cutoff no \(T>0\) works (Fable referee). The resonant value \(K=1\) used
for the \(3\kappa/4\) example requires \(\Lambda\ll\Lambda_*\); at
\(\Lambda\sim1/\tau\), \(K\sim1/(\pi\tau\omega_0)\) and the gains must
exceed \(2K\).

Small independent Gaussian readout noises of standard deviations
\(r_1,r_2\) add at most \(r_1/G\) and \(Gr_2/B\) to the two budgets
(16). Every chosen strict violation therefore permits positive readout
noise. A prescribed minimum readout noise, coupling duration, gain
ceiling or memory disturbance could restrict the protocol; those would
be added instrument premises.

\textbf{Order and cost of the limits.} For a target area, first fix a
finite cutoff and its actual prior \(V\), then choose finite gains using
(11), then choose a positive duration using (17). The constants depend
on the cutoff and the gains. Removing the cutoff before this
construction reintroduces the divergent full momentum variance. The
impulsive limit also spends unbounded control strength and allows
idealized bilinear couplings. This establishes the linear-model
obstruction; relativistic locality, material bounds and an
implementation using only electromagnetic controls would require further
estimates with \(c\). Heavy or detuned pointers are unnecessary for the
theorem. Increasing mass can reduce coordinate noise or slow damping,
but leaves the equilibrium resonant area \(\kappa\) unchanged.

\hypertarget{a-recording-law-that-restores-closure}{%
\subsection{4. A recording law that restores
closure}\label{a-recording-law-that-restores-closure}}

There is a precise point where the two-pulse protocol can be blocked.
After the first pulse, reading \(Y+GQ\) and keeping pointer 1 available
for the second pulse must itself have a conjugate cost.

\textbf{An explicit repair.} In the product resonant example, let the
stored first reading have independent Gaussian error \(N\) of variance
\(s^2>0\), and let creating that reading give pointer 1 a further
independent momentum kick \(D\) of variance \(d^2\). The body has
already received \(-G\Pi\); the second pointer now measures \(\Pi+D\).
Assume

\[s^2d^2\ge\kappa^2. \tag{20}\]

Grant arbitrarily accurate inference of \(\Pi+D\) from pointer 2. The
two independent sectors then give

\[\operatorname{Var}(Q\mid R)
=\frac{\kappa(\kappa+s^2)}{\kappa+s^2+G^2\kappa},\qquad
\operatorname{Var}(P-G\Pi\mid R)
=\kappa+\frac{G^2\kappa d^2}{\kappa+d^2}. \tag{21}\]

For \(d^2=\kappa^2/s^2\) their product is exactly \(\kappa^2\); larger
\(d^2\) or finite precision of the second read increases it. The extra
kick makes the later pointer momentum an imperfect record of the earlier
body kick. A weaker statement that every read imparts some disturbance
would leave (20) unproved.

The following sufficient premise handles arbitrary correlations and
composition. Let \(\Omega_n\) be the canonical Poisson matrix and impose
on every Gaussian joint state, including reusable apparatus and memory,

\[V+i\kappa\Omega_n\succeq0. \tag{22}\]

Admit symplectic transformations, independent ancillary states obeying
(22), discarding subsystems, and Gaussian measurements on a subsystem
that is then discarded. Such a measurement has independent Gaussian seed
covariance \(\Gamma\) with \(\Gamma+i\kappa\Omega_A\succeq0\). Sharp
single-coordinate measurements are obtained as squeezed limits. When a
pointer is retained after being read, require a dilation of its reading
within this same class; (20) is the basic coordinate example.

\textbf{Proposition 3 (Gaussian closure).} These operations preserve
(22) on the unmeasured systems conditional on the complete record,
including sequential adaptive uses. In particular the body's posterior
has \(\sqrt{\det\Sigma}\ge\kappa\).

\emph{Proof.} Symplectic transformations preserve (22), products give a
direct sum, and discarding takes a principal submatrix. For a joint
covariance
\(V=\left(\begin{smallmatrix}A&C\\C^{\sf T}&B\end{smallmatrix}\right)\),
Gaussian conditioning gives \(A-C(B+\Gamma)^{-1}C^{\sf T}\). Complex
conjugation of the seed condition gives
\(\Gamma-i\kappa\Omega_A\succeq0\). Adding that positive matrix on the
apparatus block of (22) gives

\[\begin{pmatrix}A+i\kappa\Omega_S&C\\
C^{\sf T}&B+\Gamma\end{pmatrix}\succeq0.\]

Its Schur complement proves the desired conditional inequality. Limits
preserve positivity. Apply this at every branch of a sequence, retaining
all apparatus needed at later stages. \(\square\)

This is the Gaussian quantum instrument restriction written in classical
covariance language. Bartlett, Rudolph and Spekkens derive Gaussian
quantum preparations, transformations and measurements from an epistemic
restriction (\href{https://arxiv.org/abs/1111.5057}{2012 paper}, author
abstract checked). Proposition 3 supplies the needed covariance proof
here. The common-field covariance (8) fixes the prior, which already
obeys (22); the seed restriction and the retained-pointer law are
separate premises. Law (20) is the uncorrelated case; the dilation gives
in general \(\det{\rm Cov}(N,D)\ge\kappa^2\). Even separate one-pair
bounds on every prior leave the joint condition (22) to be established.

\hypertarget{constants-complementarity-and-the-newton-comparison}{%
\subsection{5. Constants, complementarity and the Newton
comparison}\label{constants-complementarity-and-the-newton-comparison}}

The positive constant that a closure law would preserve is the resonant
\(\zeta=\kappa\). Under the additional SED/Planck identification in the
\href{sed-zeta-radiation-link.md}{revised link}, §3 (local full-read),

\[\kappa=\frac{\hbar}{2}=\frac{h_P}{4\pi},\qquad
h_*=2\zeta=2\kappa=\hbar
=\frac{(\pi^4/15)^{1/3}}{2\pi}\,h_{\rm rad}. \tag{23}\]

That thermal identification retains the conditions and disputed source
status stated there. Theorem U supplies the radiation unit independently
of any record restriction. Theorem B\('\) in the
\href{principia-fifth-postulate.md}{fifth-postulate note}, §4 (local
full-read), classifies symplectically invariant, noise-closed one-pair
Gaussian classes. Theorems 1 and 2 show why closure under
\emph{conditioning on retained records} is an additional requirement
beyond that classification.

For Proposition 3's restricted Gaussian instruments and the affine
generators of \href{newton-indeterminacy-routes.md}{Theorem F}, §4
(local full-read), Lemma G supplies statistical speed constant
\(h_*=2\kappa\). Subject to Theorem F's remaining comparison hypotheses,

\[\frac s8\sum_j\sup\Delta D_j+\frac J2\sum_j\sup\Delta X_j
\ge2\kappa\arcsin(1-2\epsilon),\qquad
s=\frac{F\tau^2}{2m},\quad J=F\tau. \tag{24}\]

The passive retained readings of Theorem 1 fall outside this restricted
class. Their small posterior leaves the unconditional impulse spread
charged in (24) in place. The unrestricted, non-Gaussian comparison
requires the further operator/statistical-speed premises of Theorem C
and Theorem F respectively.

\textbf{Theorem E's window bound.} Write
\(z_\epsilon= \Phi^{-1}(1-\epsilon/2)\) for \(0<\epsilon<1/2\), where
\(\Phi\) is the standard normal distribution function. For any
conditional Gaussian, windows centred at its conditional means with
half-widths \(a=z_\epsilon\sigma_q\) and \(b=z_\epsilon\sigma_p\)
contain the respective variables with probability \(1-\epsilon\). The
diagonal estimates in Theorem 1 give

\[ab\le z_\epsilon^2\left(
\frac{\sigma_Y\sigma_P}{G}
+\frac{\sigma_Y\sigma_{Z_2}}B\right)\longrightarrow0. \tag{25}\]

The canonical scaling (7) leaves this product unchanged. Thus the
unrestricted classical record class can violate, at any supplied
\(\hbar>0\), the necessary quantum condition

\[\lambda_0\!\left(\frac{ab}{\hbar}\right)
\ge(1-2\epsilon)^2\]

of \href{principia-fifth-postulate.md}{Theorem E}, §6b (local
full-read). Equation (25) bounds the product of marginal spreads
directly, so it also covers correlated posteriors whose determinant
alone would give insufficient control of axis-aligned windows.

Under (22) with \(\hbar=2\kappa\), every Gaussian covariance has a
Gaussian quantum-state realization: a symplectic transform of thermal
oscillator states with covariance eigenvalues at least \(\hbar/2\).
Consequently Theorem E applies to those Gaussian marginals. Extending
that statement to all non-Gaussian classical states would require an
additional state/instrument principle.

Finally the \href{newton-record-parallel-move.md}{parallel-record note},
Proposition 1 (local full-read), gives the \emph{forgotten} quantum
record's momentum variance \(\hbar^2/(4\sigma^2)=\kappa^2/\sigma^2\)
under (23). Equation (20) requires this same cost when recording a
pointer that will be reused. Forgetting a reading and preserving it for
a later kick estimate impose different conditioning questions;
Proposition 3 states a law that covers both.

\hypertarget{complete-records-the-continuity-and-noise-mechanisms-compared}{%
\subsection{6. Complete records: the continuity and noise mechanisms
compared}\label{complete-records-the-continuity-and-noise-mechanisms-compared}}

This checkpoint uses the Gaussian instrument class of Proposition 3, the
fixed-cutoff shared-bath model of Theorems 1--2, and
\href{leibniz-continuity-records.md\#1-the-proposition}{Proposition L}
(local passage). Its new statements are exact finite-dimensional
Gaussian calculations. Positivity, composition and radiation calibration
are tested separately; no additional physical law is adopted.

\hypertarget{continuity-must-concern-the-complete-apparatus-family}{%
\subsubsection{6.1 Continuity must concern the complete apparatus
family}\label{continuity-must-concern-the-complete-apparatus-family}}

Fix one body, \(m>0\), a window \([0,\tau]\) and a family
\(\mathfrak A\) of finite, non-adaptive affine apparatus protocols with
Gaussian preparations and noises independent of the force hypothesis.
Each protocol includes its entire retained memory, bath monitors, reused
pointers and terminal readings in one record \(R_\pi\). Linear dynamics
gives \[R_\pi=\mu_\pi+B_\pi\theta+Fv_\pi+\xi_\pi,
\qquad\xi_\pi\sim N(0,V_\pi),\tag{26}\] where \(\theta\) comprises the
unknown initial position and velocity. For preparation-ignorant tests
take a full maximal invariant \(Q_\pi R_\pi\), with \(Q_\pi B_\pi=0\).
Put \(b_\pi=Q_\pi v_\pi\), \(W_\pi=Q_\pi V_\pi Q_\pi^{\sf T}\) and
\[I_\pi=b_\pi^{\sf T}W_\pi^\dagger b_\pi
\quad\text{if }b_\pi\in\operatorname{Ran}W_\pi;
\qquad I_\pi=\infty\text{ otherwise}.\tag{27}\] The finite value is the
force Fisher information. The singular case has a noiseless
force-sensitive direction. The projection retains all invariant
information, including correlations with side records; it is not a
choice to discard inconvenient memory. With known preparation take
\(Q_\pi=I\) instead.

\textbf{Proposition 4 (exact uniform-continuity criterion).} In this
family, for equal prior probabilities the best invariant error is
\[P_\pi(F)=\Phi(-|F|\sqrt{I_\pi}/2).
\tag{28}\] For \(F\ne0\) the \(I_\pi=\infty\) value is zero; at \(F=0\)
it is \(1/2\). Writing \(I_*=\sup_{\pi\in\mathfrak A}I_\pi\), the
complete-record maximal-invariant laws \(P_{\pi,F}\) are uniformly
continuous at zero in total variation exactly when \(I_*<\infty\). The
quantifier covers every admitted schedule, number of marks and retained
memory: \[\sup_{\pi\in\mathfrak A}\|P_{\pi,F}-P_{\pi,0}\|_{\rm TV}
=2\Phi(|F|\sqrt{I_*}/2)-1.\tag{29}\] For \(I_*=\infty\) the right side
is one at every \(F\ne0\).

\emph{Proof.} On the supported subspace whiten the Gaussian noise. The
likelihood ratio depends only on the projection along the whitened mean
difference, whose length is \(|F|\sqrt{I_\pi}\). Thresholding at its
midpoint gives (28) and the total variation in (29). Different affine
supports in the singular case are disjoint. Monotonicity and continuity
of \(\Phi\) permit the supremum, whether or not attained. Consequently
finite \(I_*\) gives the common modulus \(|F|\sqrt{I_*}/\sqrt{2\pi}\),
while infinite \(I_*\) gives a jump. \(\square\)

For the unrestricted independent-mark family of Proposition L,
\(I_*=\tau^3/(24m\kappa)\) at \(\kappa>0\). Its zero branch has
\(I_*=\infty\). Thus that proposition uses \textbf{uniform continuity
over the family}, equivalently continuity of its optimal verdict, rather
than pointwise continuity of finite noisy apparatus. A bounded classical
family may have \(I_*<\infty\) at \(\kappa=0\). For example three
passive, independent position readings of variance \(\sigma^2\) at
\(0,\tau/2,\tau\), with zero recoil, give the invariant statistic
\(R_0-2R_{\tau/2}+R_\tau\), mean \(F\tau^2/(4m)\) and variance
\(6\sigma^2\). Its information is \[I=\frac{\tau^4}{96m^2\sigma^2}.\]
More generally a finite total precision budget
\(\sum_j\delta_j^{-2}\le C\) gives \(I_*\le C\tau^4/(4m^2)\) even with
known preparation. Such budgets add under composition; maintaining the
same \(C\) requires a resource restriction. Unlimited repeats on the
same zero-recoil trajectory make \(I_*\to\infty\).

Therefore (29) proves the exact statistical requirement of the
continuity route. Supplying it independently for a physically specified
family closed under arbitrary finite composition remains necessary. It
alone supplies neither a terminal imprecision--kick law nor an action
constant common to different \(m,\tau\) and apparatus classes. The
quantity \(\tau^3/(24mI_*)\) has action units when \(0<I_*<\infty\);
identifying it with a universal mark scale uses the additional
mark-model saturation and universality premises.

\hypertarget{the-noise-law-must-survive-side-observations-of-the-kick}{%
\subsubsection{6.2 The noise law must survive side observations of the
kick}\label{the-noise-law-must-survive-side-observations-of-the-kick}}

For any jointly Gaussian body \(X\) and complete record \(\mathcal R\),
the exact answer is the supported Schur complement
\[\operatorname{Cov}(X\mid\mathcal R)
=K_{XX}-K_{X\mathcal R}K_{\mathcal R\mathcal R}^\dagger
 K_{\mathcal R X}.\tag{30}\] Appending a record decreases this
covariance in the positive-semidefinite order. A prior noise covariance
therefore needs a conditioning estimate, including indirect
observations, before it can enforce a floor.

Here is an exact extension of (20)--(21). Start with the algebraic
resonant example \(V=\kappa I_6\), \(\kappa>0\), in canonical scaled
coordinates. The first pulse still gives body \(X=(Q,P-G\Pi)\) and
pointer coordinate \(Y+GQ\). Storing the latter produces error \(N\) and
kicks the reusable pointer by \(D\). Grant exact inference of its later
momentum \(K=\Pi+D\) from the second pointer, as in §4's limiting second
readout. Let \((N,D,S)\) be jointly Gaussian and independent of the six
initial body--pointer variables; \(S\) is any side record of the readout
noise. It can be a monitor read earlier or later, or stored in a
separate memory. The complete record is \[\mathcal R=(Y+GQ+N,\ K,\ S).\]
Subtract conditional means given \(S\) and put \[\begin{aligned}
E_S&=\operatorname{Cov}\bigl((N,D)^{\sf T}\mid S\bigr)
=\begin{pmatrix}a&c\!_{ND}\\c\!_{ND}&b\end{pmatrix}\succeq0,\\
x&=G^2,\quad u=\kappa+a,\quad v=\kappa+b,\\
\mathcal D&=(\kappa x+u)v-c\!_{ND}^{\,2}>0.
\end{aligned}\tag{31}\] The symbol \(c\!_{ND}\) is a covariance,
distinct from the speed \(c\). All of \(a,b,c\!_{ND},u,v\) have action
units.

\textbf{Proposition 5 (the conditional recording law, necessary and
sufficient in this model).} The exact posterior determinant is
\[\boxed{\det\operatorname{Cov}(X\mid\mathcal R)
=\kappa^2\left[1+\frac{x\bigl(ab-c\!_{ND}^{\,2}-\kappa^2\bigr)}
                         {\mathcal D}\right].}\tag{32}\] For every
nonzero gain \(G\), the posterior area is at least \(\kappa\) if and
only if \(\det E_S\ge\kappa^2\). A finite-precision second reading can
only increase the posterior covariance; the exact-reading limit gives
the sharp necessity test.

\emph{Proof.} Condition first on \(S\). The remaining two readings have
covariance
\(B_R=\left(\begin{smallmatrix}\kappa x+u&c\!_{ND}\\ c\!_{ND}&v\end{smallmatrix}\right)\),
while \(A_X=\operatorname{diag}(\kappa,\kappa(1+x))\) and
\(C_{XR}=\operatorname{diag}(G\kappa,-G\kappa)\). In particular (30)
gives \[\begin{aligned}
\Sigma_{11}&=\kappa(uv-c\!_{ND}^{\,2})/\mathcal D,\\
\Sigma_{12}&=-x\kappa^2c\!_{ND}/\mathcal D,\\
\Sigma_{22}&=\kappa(1+x)-x\kappa^2(\kappa x+u)/\mathcal D.
\end{aligned}\tag{33}\] For a compact determinant check, the covariance
of the readings conditional on \(X\) is
\(\left(\begin{smallmatrix}u&c\!_{ND}\\c\!_{ND}&b+\kappa/(1+x) \end{smallmatrix}\right)\).
Taking the joint determinant in both block orders yields \[\det\Sigma
=\kappa^2\frac{\kappa u+(1+x)(ub-c\!_{ND}^{\,2})}
                    {(\kappa x+u)v-c\!_{ND}^{\,2}}.\] The numerator
minus the denominator is \(x(ab-c\!_{ND}^{\,2}-\kappa^2)\), proving
(32). The initial strictly positive \(\kappa I_6\) makes the reading
covariance invertible, including singular \(E_S\). \(\square\)

\textbf{A concrete indirect-monitor test.} Let \(N,D\) be independent
with \(\operatorname{Var}N=a_0>0\), \(\operatorname{Var}D=b_0>0\) and
\(a_0b_0=\kappa^2\), so the marginal recording law (20) is saturated.
Allow the side record \(S=D+M\), where \(M\) is an independent Gaussian
monitor error of variance \(r^2>0\). Then
\[a=a_0,\qquad b=\frac{b_0r^2}{b_0+r^2},\qquad c\!_{ND}=0,
\qquad\det E_S=\kappa^2\frac{r^2}{b_0+r^2}<\kappa^2.\tag{34}\] Every
informative monitor, even a noisy one, makes (32) strictly smaller than
\(\kappa^2\) for \(G\ne0\). At \(r\downarrow0\), the kick is known and
\(\det\Sigma=\kappa^2u/(u+\kappa G^2)\to0\) as \(G\to\infty\). At
\(r\to\infty\) the original saturated repair is recovered. All limits
keep their order and apparatus assumptions explicit. A realizable
monitor may itself disturb the pointer; the passive monitor in this test
is a granted operation whose physical admissibility is exactly what a
closure mechanism must settle.

For arbitrary correlated readout noise with covariance \(E_0\) and side
record covariance \(V_S\), the matrix to test is
\[E_S=E_0-\operatorname{Cov}((N,D),S)V_S^\dagger
                 \operatorname{Cov}(S,(N,D)).\tag{35}\] The common
bath's spectrum can set \(E_0\), but it leaves this conditional
innovation to be controlled. In the general correlated prior of (8), use
(30) on the full joint record; (32)'s iff is asserted for its stated
product-prior example only. Theorems 1--2 already cover arbitrary finite
body--pointer cross correlations. Proposition 3 supplies a sufficient
composition law on every retained subsystem; (34) identifies a side
observation that must be priced or excluded by that law.

\hypertarget{ideas-supplied-by-the-current-lean-formalization}{%
\subsubsection{6.3 Ideas supplied by the current Lean
formalization}\label{ideas-supplied-by-the-current-lean-formalization}}

The sister repository is read at commit \texttt{381d5e9}, keeping its
historical theorems and modern diagnostics separate. Its
\href{https://github.com/arivero/newtonlean/blob/381d5e9/research/graphs.md}{dependency
graph} and
\href{https://github.com/arivero/newtonlean/blob/381d5e9/research/PROP_I_REALIZATION.md}{Proposition
I realization obligations} (local full-read) expose a specific missing
bridge in both printed stages: the cited Lemma III Cor. 4 concerns
approximating a given curve, while the force-generated vertices move
under refinement. The finite area law (P1) and refinement identities
(P2) are checked; realization (P3) and force identification (P5) remain
separate. The harmonic field has a mesh-uniform quadratic stability
invariant. The pending norm inequality in its Fraction representation is
a formalization obstacle.
\href{refinement-composition-and-limit.md\#3b-harmonic-polygons-a-finite-refinement-bound-from-a-different-norm}{The
refinement note, §3b} now proves the finite stability/defect argument in
a weighted sum norm, including unequal partitions and partial cells. It
is a written result, separate from the sister repository's formalized
statements, and supplies no action floor.

The graph suggests the same distinction for records: geometric control
of the body and apparatus map needs an additional statistical control in
the complete-record norm. The following finite declarations were read in
their definitions and proofs (local full-read). Their algebra supplies
two concrete tests for that statistical bridge. The Gaussian
consequences here are written derivations, separate from the compiled
finite theorems.

First,
\href{https://github.com/arivero/newtonlean/blob/381d5e9/NewtonLimitDynamics/Polygon/RelativeMotion.lean}{RelativeMotion.corVI\_relative
and relative\_deflection\_difference} derive exact cancellation of an
arbitrary common deflection history from both relative coordinates. A
classical Gaussian consequence is immediate: if two observed motions
receive the same acceleration history, differencing them removes that
random history too. More generally if a retained signal is \(r+L\eta\),
with bath covariance \(C_\eta\), a difference map \(D_r\) leaves
covariance \[D_rL C_\eta L^{\sf T}D_r^{\sf T}.\tag{36}\] If \(D_rL=0\)
it is zero regardless of the bath amplitude. This is a modern
noise-model consequence of the finite cancellation identity, not a
historical assertion about radiation. A noise floor must survive
accessible relative observations as well as the conditional subtraction
(35). Common driving alone does not guarantee that survival.

Second,
\href{https://github.com/arivero/newtonlean/blob/381d5e9/NewtonLimitDynamics/Polygon/PartitionComparison.lean}{PartitionComparison.partition\_gap
and PartialCell.candidate\_partial\_residual} and
\href{https://github.com/arivero/newtonlean/blob/381d5e9/NewtonLimitDynamics/Polygon/PartialCell.lean}{the
partial-cell proof} separate the actual end-kick polygon from its
rational constant-force comparison map. At any rational sample time
\(t\le\tau\) the exact position residual is \(a_{\rm acc}c_\pi(t)\),
where \[c_\pi(t)=\tfrac12\left(\sum_{j\text{ complete}}h_j^2+u^2\right),
\qquad0\le c_\pi(t)\le\tfrac12t|\pi|,\qquad
a_{\rm acc}=F/m.\tag{37}\] Here \(u\) is the partial final duration and
\(|\pi|\) bounds every complete and partial cell. The formal theorem
uses represented fractions and cross-multiplication equivalence; (37) is
its rational-value reading. It fixes the force signal's geometric
discrepancy uniformly over partitions. It supplies no measurement law.

For a single Gaussian position record of width \(\sigma_\pi\) the total
variation between polygon and corrected-map records is exactly
\[2\Phi\left(\frac{|F|c_\pi(t)}{2m\sigma_\pi}\right)-1.\tag{38}\] Thus
the mechanically vanishing residual gives statistically convergent
records only when its noise-normalized size vanishes. A width
\(\sigma_\pi=F_0c_\pi(t)/m\) with fixed force \(F_0>0\) leaves (38)
constant for every nonzero \(F\), even as the mesh vanishes. A fixed
positive width makes it tend to zero. For the complete multivariate
record the corresponding condition is
\(\|W_\pi^{-1/2}Q_\pi(v_\pi-v_{\rm cont})\|\to0\) on its supported
subspace; a noiseless residual direction again has total variation one.
This distinguishes uniform mechanical refinement from uniform
complete-record refinement by an explicit calculation.

The same repository's
\href{https://github.com/arivero/newtonlean/blob/381d5e9/NewtonLimitDynamics/Diagnostic/PhaseArea.lean}{PhaseArea.cells\_area2}
preserves every phase triangle under affine drift/kick schedules. It
supports the action \emph{kind} of a symplectic covariance restriction,
while preserving triangles of every initial size. Its invariant gives no
selected minimum. The end-kick convention of (37) differs from the
symmetric half-kick construction of the
\href{refinement-composition-and-limit.md}{insertion note}; each
statement retains its own convention. These tests concern records and
the polygon--trajectory discrepancy, rather than Kepler sectors.

\hypertarget{mechanism-selected-and-the-next-lemma}{%
\subsubsection{6.4 Mechanism selected and the next
lemma}\label{mechanism-selected-and-the-next-lemma}}

The noise route now has a sharper constructive target: a body--apparatus
model must retain a \textbf{conditional innovation pair} with
determinant at least \(\kappa^2\) after every admissible side record.
Its readout dynamics must charge observations of the kick and of
reusable memory, as (34) shows. The immediate lemma is a branchwise
innovation estimate for finite adaptive linear recording, followed by
independence from the refinement schedule. Reject a proposed common-bath
mechanism if an admissible relative channel annihilates its noise as in
(36), or if a terminal side record reduces its innovation determinant
below \(\kappa^2\) as in (34).

The continuity route has been reduced to (29), but no independently
derived uniform bound over a composition-closed physical family is
available. Neither route has discharged positivity from independent
physical premises. We therefore select the conditional-innovation
calculation as this checkpoint's advance, and the innovation law as the
next mechanism to construct. No new physical premise has been installed.
Within an admitted Gaussian closure law \(h_*=2\kappa>0\); proving that
law, proving its positive scale and identifying
\(h_*=\gamma h_{\rm rad}\) remain distinct. Formula (23)'s radiation
calibration retains its existing conditional source status.

\hypertarget{adaptive-records-innovation-replacement-and-refinement}{%
\subsection{7. Adaptive records, innovation replacement and
refinement}\label{adaptive-records-innovation-replacement-and-refinement}}

The exact Gaussian kernels below turn the complete-record requirement
into a composition calculation. The harmonic estimate developed from the
sister formalization supplies a finite mechanical input. All statements
retain their declared apparatus and covariance assumptions; no
electromagnetic readout law is inferred from them.

\hypertarget{an-exact-likelihood-identity-for-adaptive-recording}{%
\subsubsection{7.1 An exact likelihood identity for adaptive
recording}\label{an-exact-likelihood-identity-for-adaptive-recording}}

Fix a protocol with finitely many readings \(R_1,\ldots,R_n\), each a
finite-dimensional real vector. Include all random seeds, retained
memory, noise monitors and terminal readings. Conditional on the
previous complete history \(r_{<j}\), suppose \[R_j\mid(r_{<j},F)\sim
N\bigl(\mu_j(r_{<j})+F b_j(r_{<j}),\ W_j(r_{<j})\bigr),
\qquad W_j\succ0.\tag{39}\] Independent randomized controls may use any
force-independent law; retain their values in the initial history, so
their common density cancels in (41). Count identical copies of an
already stored reading once; deterministic memory functions remain
recoverable from the complete history. The protocol may choose its next
gains, time and apparatus from that history. At a fixed history those
choices are the same under both force hypotheses, and \(W_j,\mu_j,b_j\)
are independent of \(F\). For a linear Gaussian body--apparatus model
with a declared Gaussian preparation, this follows by successive
Gaussian conditioning: once the realized gains are fixed, the
conditional mean is affine in \(F\) and the covariance is independent of
\(F\). The complete joint record can be non-Gaussian because the gains
depend on previous readings. Preparation-ignorant adaptive tests require
their own justified invariant kernels; §6.1's non-adaptive projection
alone does not establish (39) for them.

Define the nonnegative information along a branch by
\[I_j(r_{<j})=b_j^{\sf T}W_j^{-1}b_j,
\qquad J(r)=\sum_{j=1}^n I_j(r_{<j}).\tag{40}\] Both \(J\) and its cap
\(C\) have inverse force squared units, so \(F^2J\) is dimensionless.
Let \(p_F(r)\) be the normalized complete-record density obtained by
multiplying the conditional kernels. Its Hellinger affinity with \(p_0\)
is \(\mathcal A(F,0)=\int\sqrt{p_Fp_0}\,dr\).

\textbf{Proposition 6 (adaptive affinity and continuity).} Exactly,
\[\sqrt{p_F(r)p_0(r)}
=p_{F/2}(r)e^{-F^2J(r)/8},\qquad
\mathcal A(F,0)=E_{F/2}e^{-F^2J(R)/8}.\tag{41}\] For a family of these
protocols, uniform complete-record continuity at \(F=0\) is equivalent
to \[\sup_{\pi}E_{\pi,F/2}
       \left[1-e^{-F^2J_\pi(R)/8}\right]\longrightarrow0.
\tag{42}\] A sufficient condition, uniform in every admitted schedule,
history and retained memory, is \(J_\pi(r)\le C<\infty\). It gives
\[\|P_{\pi,F}-P_{\pi,0}\|_{\rm TV}
\le\sqrt{1-e^{-F^2C/4}}\le\tfrac12|F|\sqrt C.\tag{43}\] This
deterministic budget is sufficient; rare high-information branches need
only satisfy the exact averaged condition (42).

\emph{Proof.} At the same history, completing the Gaussian square gives
\[\sqrt{g_j(r_j;F)g_j(r_j;0)}
=g_j(r_j;F/2)e^{-F^2 I_j(r_{<j})/8}.\] Multiply these identities, then
integrate, proving (41). For any two normalized densities with affinity
\(\mathcal A\), \[1-\mathcal A\le\|P-Q\|_{\rm TV}
\le\sqrt{1-\mathcal A^2}.\] The lower bound follows from
\(\min(p,q)\le\sqrt{pq}\). For the upper bound apply Cauchy--Schwarz to
\(|p-q|=|\sqrt p-\sqrt q|(\sqrt p+\sqrt q)\) and divide its integral by
two. These inequalities prove (42); (41) and \(J\le C\) prove (43),
using \(1-e^{-u}\le u\). \(\square\)

Information also composes exactly. The score is the sum of
\(b_j^{\sf T}W_j^{-1}(R_j-\mu_j-Fb_j)\). Each summand has zero
conditional mean and conditional variance \(I_j\). When \(E_FJ<\infty\),
orthogonality of these martingale differences gives full-record Fisher
information \[\mathcal I(F)=E_FJ(R).\tag{44}\] Joining two blocks of
readings adds their branchwise \(J\) values, with the second block
conditioned on the complete first history. Budgets therefore add; a
fixed cap \(C\) is not preserved under unrestricted repetition merely
because each apparatus has a finite cap. For a non-adaptive Gaussian
record \(J\) is deterministic, recovering §6.1's criterion. A noiseless
force-sensitive conditional direction violates the finite-budget
hypothesis; it cannot be discarded as irrelevant memory.

\hypertarget{a-side-record-spends-information-and-removes-innovation}{%
\subsubsection{7.2 A side record spends information and removes
innovation}\label{a-side-record-spends-information-and-removes-innovation}}

An indirect observation can carry force information even when its
marginal law is force-independent. For example let \(R=Fv+N\),
\(S=N+M\), with independent centred Gaussian \(N,M\) of variances
\(n^2,r^2>0\). In record order \(R,S\), the first kernel has
\(I_1=v^2/n^2\), and \(S\mid(R,F)\sim N(R-Fv,r^2)\) has \(I_2=v^2/r^2\).
Consequently \[J=v^2(n^{-2}+r^{-2}).\tag{45}\] This agrees with
conditioning the noise on \(S\), whose variance is \(n^2r^2/(n^2+r^2)\).
The raw fact that \(S\) has no force-dependent mean does not make it
harmless after retaining \(R\).

There is also a quantitative repair of §6.2's two-pointer monitor test.
Start with its saturated independent error and kick,
\(a_0b_0=\kappa^2\), and side reading \(S=D+M\) with independent monitor
error \(M\) of variance \(r^2\). After this monitor, apply a fresh
centred Gaussian kick \(E\) to the reusable pointer momentum,
independent of the initial variables, \(N,D,S\) and of variance \(e^2\).
The second pointer now infers \(K=\Pi+D+E\). By Proposition 5, for a
nonzero first gain the body's conditional area remains at least
\(\kappa\) precisely when
\[a_0\left[\frac{b_0r^2}{b_0+r^2}+e^2\right]\ge\kappa^2,
\quad\text{equivalently}\quad
e^2\ge\frac{b_0^2}{b_0+r^2}.\tag{46}\] The required fresh variance is
exactly the part of \(D\) learned by the monitor,
\(\operatorname{Var}(E[D\mid S])\). At an uninformative monitor
\(r\to\infty\) its cost tends to zero; at an exact monitor
\(r\downarrow0\) it must replace all \(b_0\). Equality restores the
determinant to \(\kappa^2\) for every gain in the stated product-prior
model. All these variances have action units in the canonical scaling;
(46)'s product has action squared units.

This is a useful conditional construction, with the extra kick openly
supplied. A later admissible side record of \(E\) must be included and
would reduce its innovation again. A monitor dynamics that necessarily
creates inaccessible fresh noise of at least (46), including its own
retained memory, is the physical mechanism to derive. Adding noise by
prescription proves neither that mechanism nor \(\kappa>0\).

\hypertarget{mechanical-refinement-in-the-complete-conditional-record-norm}{%
\subsubsection{7.3 Mechanical refinement in the complete conditional
record
norm}\label{mechanical-refinement-in-the-complete-conditional-record-norm}}

For two record constructions at the same force, suppose their
conditional kernels have the same \(W_j(r_{<j})\succ0\) and means
\(m_{\pi,j}(r_{<j})\), \(m_{\sigma,j}(r_{<j})\). Put
\(\delta_j=m_{\pi,j}-m_{\sigma,j}\) and
\[D_{\pi\sigma}(r)=\sum_j\delta_j^{\sf T}W_j^{-1}\delta_j.
\tag{47}\] Completing the same square as in (41) gives affinity
\(E_{\rm mid}e^{-D_{\pi\sigma}/8}\), where the normalized midpoint law
uses the average conditional means and the common conditional
covariances. In particular, if every branch has
\(D_{\pi\sigma}\le\epsilon_{\pi\sigma}^2\), then
\[\|P_\pi-P_\sigma\|_{\rm TV}
\le\sqrt{1-e^{-\epsilon_{\pi\sigma}^2/4}}
\le\epsilon_{\pi\sigma}/2.\tag{48}\] This is an adaptive complete-record
refinement criterion. Covariances that change under refinement need an
additional comparison; equality of the kernels' covariance is a stated
hypothesis here.

For an explicit mechanical input use the harmonic end-kick result in
\href{refinement-composition-and-limit.md\#3b-harmonic-polygons-a-finite-refinement-bound-from-a-different-norm}{the
refinement note, Proposition 3b} (local full-read). Fix
\(z_0,w,\rho,\tau\) as there. Suppose at every common history the
conditional mean discrepancy is
\(\delta_j=H_j(r_{<j})[z_\pi(t_j)-z_\sigma(t_j)]\), with the same
adaptive sampling time \(t_j\) under both constructions. Define
\(\ell_j=\|W_j^{-1/2}H_j\|_{\rho\to2}\) and require the branchwise
budget \(\sum_j\ell_j^2\le L^2\), uniform in schedules and memory.
Equations (6f) there and (47)--(48) then prove
\[\|P_\pi-P_\sigma\|_{\rm TV}
\le\frac L2\,w e^{\rho\tau}(\tau+\rho^{-1})
 (|\pi|+|\sigma|)\|z_0\|_\rho.\tag{49}\] Every constant is explicit. The
estimate permits arbitrarily many readings when their total whitened
sensitivity obeys the declared budget. If the noise widths shrink so
fast that this budget diverges, the mechanical Cauchy bound supplies no
statistical convergence, as the exact constant-force example (38)
already demonstrates. The same-noise linear readout model gives a
concrete instance; a general back-reacting body--apparatus model must
prove the assumed conditional mean comparison and covariance control.

\hypertarget{next-mechanism-and-abandonment-test}{%
\subsubsection{7.4 Next mechanism and abandonment
test}\label{next-mechanism-and-abandonment-test}}

Adaptivity no longer prevents an exact finite likelihood calculation:
(41) includes its realized gains and complete retained memory, and (49)
supplies a sufficient schedule-independent record comparison. The
missing physical input is now concrete. Construct a coupled
body--pointer--monitor--memory dynamics in which observing the pointer
kick necessarily replaces the learned variance by fresh conditional
innovation, with (46) as the one-monitor test. It must survive a
terminal read of that memory and the accessible relative channels in
(36). Abandon the mechanism if either operation removes its innovation;
a marginal bath spectrum cannot repair that failure.

The budget in (43) can hold for a bounded classical apparatus at the
zero-action branch, while unrestricted zero-recoil repeats make it
diverge. The repair in (46) is proportional to a supplied \(\kappa^2\)
through \(a_0b_0=\kappa^2\) and becomes vacuous on the zero branch.
Hence this checkpoint establishes adaptive composition and a finite
refinement estimate under explicit assumptions, and quantifies the
required innovation replacement. It does not derive positivity,
universality or the radiation calibration in (23).

\hypertarget{hamiltonian-memory-terminal-records-and-blocking}{%
\subsection{8. Hamiltonian memory, terminal records and
blocking}\label{hamiltonian-memory-terminal-records-and-blocking}}

This is a conditional finite body--apparatus--record construction.
Assume the initial Gaussian canonical state satisfies (22) with
\(\kappa>0\), allow affine symplectic evolution, and realize the record
as physical memory coordinates read at the terminal cut. These are
independently stated model hypotheses, rather than a separately
postulated measurement seed. Their general physical necessity and the
initial positive scale are not established by this construction.

\hypertarget{a-monitor-that-disturbs-the-physical-stored-reading}{%
\subsubsection{8.1 A monitor that disturbs the physical stored
reading}\label{a-monitor-that-disturbs-the-physical-stored-reading}}

Retain §6.2's product resonant body--pointer preparation and first
pulse. The body is \(X=(Q,P-G\Pi_1)\) and pointer coordinate is
\(Y_1'=Y_1+GQ\). Introduce independent memory \(A=(Z_A,\Lambda_A)\) with
diagonal Gaussian variances \(a_0,b_0>0\) and \(a_0b_0=\kappa^2\). A
unit pulse with integrated Hamiltonian generator \(Y_1'\Lambda_A\)
copies the pointer into its coordinate:
\[R_1=Z_A+Y_1',\qquad K=\Pi_1-\Lambda_A.
\tag{50}\] All expressions here use the initial random variables. The
body stays unchanged: \(\{P-G\Pi_1,Y_1'\}=-G+G=0\). Thus the first
imprecision and reusable-pointer kick are \(N=Z_A\), \(D=-\Lambda_A\),
satisfying (20) mechanically from the memory's prior. The copying pulse
is a symplectic map, not a passive assignment of an externally stored
number.

Now introduce independent monitor memory \(B=(Z_B,\Lambda_B)\) with
diagonal variances \(r^2,d^2>0\) and \(r^2d^2\ge\kappa^2\). A unit pulse
generated by \(\Lambda_A\Lambda_B\) changes both memory coordinates. The
complete terminal readings are \[R_1'=Z_A+Y_1'+\Lambda_B,\qquad
S=Z_B+\Lambda_A.\tag{51}\] The pointer momentum \(K\) is unchanged;
grant the same exact second-pointer inference of \(K\) as in Proposition
5. The effective noise is \(N'=Z_A+\Lambda_B\), \(D=-\Lambda_A\).
Conditional on \(S\), \[\begin{aligned}
a&=a_0+d^2,\qquad b=\frac{b_0r^2}{b_0+r^2},\qquad c\!_{ND}=0,\\
\det E_S&=\kappa^2+
          \frac{b_0(r^2d^2-\kappa^2)}{b_0+r^2}\ge\kappa^2.
\end{aligned}\tag{52}\] The diagonal product preparation makes \(N'\)
independent of \(D,S\); the remaining variance is the scalar Gaussian
conditioning formula. Equality holds when the monitor memory is also
saturated, \(r^2d^2=\kappa^2\). Equation (32) then gives body posterior
area exactly \(\kappa\) for every first gain. Finite second-pointer
precision increases it. A sharper monitor \(r\downarrow0\) requires
\(d^2\ge\kappa^2/r^2\) and increasingly disturbs the stored first
coordinate. At an uninformative monitor \(r\to\infty\) with saturated
\(d^2\), its extra disturbance tends to zero. These are finite-model
limits; no cutoff removal is implied.

The causal distinction from (34) is measurable in the Poisson brackets:
\[\{R_1',S\}=1-1=0,\qquad
\{R_1,S\}=1,\qquad\{R_1',K\}=1-1=0.\tag{53}\] The old numerical value
\(R_1\) and the later monitor reading are not the two terminal
coordinates of this closed memory construction. If one grants a passive
external read of \(R_1\) before the monitor and retains that number
without another physical memory coupling, the old escape of (34)
returns. Copying it to another physical memory adds that memory's
variables and conjugate kick to the calculation. The monitor's
disturbance prices an observation of the earlier kick by changing the
stored imprecision; it differs from §7.2's prescribed fresh momentum
kick but satisfies the same conditional determinant test. All variances
in (52) have action units in the canonical scaling.

\hypertarget{a-general-posterior-theorem-for-complete-terminal-coordinates}{%
\subsubsection{8.2 A general posterior theorem for complete terminal
coordinates}\label{a-general-posterior-theorem-for-complete-terminal-coordinates}}

Let \(Z\) be a centred finite jointly Gaussian canonical vector with
covariance \(V\) and Poisson matrix \(\Omega\), satisfying
\(V+i\kappa\Omega\succeq0\). Write the remaining canonical body and
complete record as \(X=LZ\), \(R=CZ\), omitting irrelevant affine means.
Suppose \[L\Omega L^{\sf T}=\Omega_X,\qquad
L\Omega C^{\sf T}=0,\qquad C\Omega C^{\sf T}=0.
\tag{54}\] These conditions hold after any affine symplectic evolution
when \(X\) is a remaining subsystem and \(R\) consists of coordinate
readings of distinct terminal memory pairs. The last condition means
their classical Poisson brackets vanish, also called an isotropic
record. This statement uses Poisson geometry; it assumes no quantum
operator commutator. General simultaneous classical records need not
satisfy (54).

\textbf{Proposition 7 (terminal-memory posterior closure).} Under these
hypotheses, \[\Sigma_{X\mid R}
=LVL^{\sf T}-LVC^{\sf T}(CVC^{\sf T})^\dagger CVL^{\sf T},
\qquad\Sigma_{X\mid R}+i\kappa\Omega_X\succeq0.
\tag{55}\] For one body pair this gives
\(\det\Sigma_{X\mid R}\ge\kappa^2\). No independent measurement seed or
readout product law is assumed.

\emph{Proof.} Push the initial positive Hermitian matrix through
\((L,C)\). By (54) its blocks are \[\begin{pmatrix}
LVL^{\sf T}+i\kappa\Omega_X&LVC^{\sf T}\\
CVL^{\sf T}&CVC^{\sf T}
\end{pmatrix}\succeq0.\] The Schur complement proves (55). For a
singular record covariance, its null directions have zero cross
covariance with \(X\); restrict to its range, or use the displayed
pseudoinverse. The one-pair Hermitian determinant is
\(\det\Sigma-\kappa^2\), so positivity gives the area bound. \(\square\)

This is an explicit Hamiltonian realization of Proposition 3's Gaussian
restriction. It reduces the readout-law obligation to a joint
preparation condition and physical storage geometry within this affine
class. It does not prove that every physically admissible observation
belongs to the class. In particular, treating a reading as an external
classical number and subsequently probing its retained conjugate
bypasses (54), exactly as (53) shows. Individual marginal phase-area
floors do not imply the initial joint condition when apparatus and field
are correlated; that full condition must be checked.

At a fixed adaptive branch the same proof applies if the full
conditional remaining state already satisfies (22), and the next map and
added Gaussian ancillas obey the stated hypotheses. It gives a
conditional induction for Hamiltonian copies and terminal reads with the
measured pair removed from subsequent mechanical use. It does not prove
closure for arbitrary nonlinear Hamiltonian feedback or for future
observations of a discarded conjugate.

\hypertarget{what-survives-blocking-and-what-must-be-proved-across-depths}{%
\subsubsection{8.3 What survives blocking, and what must be proved
across
depths}\label{what-survives-blocking-and-what-must-be-proved-across-depths}}

Fix the same terminal body variable \(X\). Let a coarse record be a
deterministic linear block of a finer terminal record, \(R_c=BR\). Its
pulled-back matrix is \(C_c=BC\), so (54) survives exactly. The coarse
body posterior also satisfies (55), and \[\operatorname{Cov}(X\mid R_c)
\succeq\operatorname{Cov}(X\mid R).\tag{56}\] To see the latter,
condition the finer posterior mean on \(R_c\): the covariance of that
mean is the nonnegative extra term in the Gaussian conditional-variance
decomposition. Repeated linear blocking composes its \(B\) matrices and
preserves this conditional floor at each finite step. This is an
explicit retained-observable statement, independent of whether a
physical trajectory curve has been constructed.

It still needs cross-depth consistency. First construct one genuine
two-storage-cell refinement and its explicit coarse record projection,
with both descriptions referring to the same terminal body. If finer
dynamics adds new copies, compare its joint body--coarse-record law with
the earlier one after that projection. The extra copy kick is part of
the comparison; forgetting a new record can disturb the body, as the
\href{newton-record-parallel-move.md}{parallel-record construction}
(local full-read) already quantifies. A sufficient next lemma is a
summable bound on that projected-law defect for a physically stated copy
schedule, while retaining (55) and useful force sensitivity. Covariance
monotonicity alone cannot compare different terminal bodies. Abandon the
schedule if matching its body requires removing an unobserved physical
kick, if its old record is preserved only by a passive external
assignment, if its complete noise can be removed by an accessible
relative or terminal-memory channel, or if its total conditional
information loses the continuity in (42).

The zero branch remains admissible: at \(\kappa=0\) the covariance
condition is ordinary positivity, and sharp preparation/copy limits
carry no universal positive cost. Thus this construction discharges a
separate seed law within the declared closed affine apparatus class,
while retaining the positive joint preparation scale as an input.
Establishing that preparation law from independent physical premises,
universality across apparatus and the calibration
\(h_*=\gamma h_{\rm rad}\) remain distinct tasks.

\hypertarget{two-physical-copies-at-distinct-times}{%
\subsubsection{8.4 Two physical copies at distinct
times}\label{two-physical-copies-at-distinct-times}}

Fix one body of mass \(m>0\) with canonical position and momentum
\((q,p)\), drifting under \(H=p^2/(2m)-Fq\) for a deterministic force
\(F\). Its centred initial Gaussian has diagonal variances \(A,B>0\),
\(AB\ge\kappa^2\). Let \(r>0\) be a supplied precision per unit time.
For \(u,v>0\), \(u+v=h\), prepare two independent memory pairs
\((Z_i,\Lambda_i)\), independent of the body, with
\(\operatorname{Var}Z_i=1/(rh_i)\),
\(\operatorname{Var}\Lambda_i=\kappa^2rh_i\), \(h_1=u\), \(h_2=v\), and
zero coordinate--momentum covariance. Their coordinates have position
units and their momenta have momentum units. This satisfies the joint
condition (22) for \(\kappa>0\). The same classical formulas also make
sense at \(\kappa=0\) with zero memory kicks.

At times \(u\) and \(h\) apply unit symplectic copy pulses generated by
\(q\Lambda_1\) and \(q\Lambda_2\), respectively. Apart from these pulses
the memories are frozen. Immediately after the second pulse, the body
and all terminal memory coordinates are exactly \[\begin{aligned}
Q_f&=q+hp/m+Fh^2/(2m)-v\Lambda_1/m,\\
P_f&=p+Fh-\Lambda_1-\Lambda_2,\\
R_1&=q+up/m+Fu^2/(2m)+Z_1,\\
R_2&=q+hp/m+Fh^2/(2m)-v\Lambda_1/m+Z_2.
\end{aligned}\tag{57}\] These are physical terminal coordinates, not
passive earlier reads. For example \(\{R_1,Q_f\}=v/m-v/m=0\) and
\(\{R_1,P_f\}=1-1=0\). Likewise
\(\{R_2,Q_f\}=\{R_2,P_f\}=\{R_1,R_2\}=0\) and \(\{Q_f,P_f\}=1\). The
later copy leaves the first stored coordinate unchanged because it
already commutes with the evolved body. Thus (54)--(55) apply to the
complete fine record \((R_1,R_2)\).

Put \(\alpha=u/h\), \(\beta=v/h\), and form the block
\(R_c=\alpha R_1+\beta R_2\). The initial memories admit the exact
canonical change of variables \[\begin{aligned}
Z_c&=\alpha Z_1+\beta Z_2,& L_c&=\Lambda_1+\Lambda_2,\\
Z_{\rm rel}&=Z_1-Z_2,&
L_{\rm rel}&=\beta\Lambda_1-\alpha\Lambda_2,\\
\bigl(\operatorname{Var}Z_c,\operatorname{Var}L_c\bigr)
 &=\bigl((rh)^{-1},\kappa^2rh\bigr),\\
\bigl(\operatorname{Var}Z_{\rm rel},
      \operatorname{Var}L_{\rm rel}\bigr)
 &=\bigl(h/(ruv),\kappa^2ruv/h\bigr).
\end{aligned}\tag{58}\] Both pairs have bracket one, all cross brackets
and cross covariances vanish, and the pairs are Gaussian independent. In
particular \(\Lambda_1=\alpha L_c+L_{\rm rel}\). At one copy time the
relative mode would decouple; the distinct times in (57) leave a real
residue.

\hypertarget{one-effective-coarse-copy-and-the-joint-law-defect}{%
\subsubsection{8.5 One effective coarse copy and the joint-law
defect}\label{one-effective-coarse-copy-and-the-joint-law-defect}}

Set \(s=h-\alpha v\in(u,h)\) and \(c=\alpha\beta v/m\). The coarse
comparator prepares the effective pair \((Z_c,L_c)\), applies the
symplectic memory shear \(Z_c\mapsto Z_c-cL_c\), copies the body
position into that memory at \(s\), and drifts the body to \(h\). Its
terminal variables are \[\begin{aligned}
Q_c&=q+hp/m+Fh^2/(2m)-\alpha v L_c/m,\\
P_c&=p+Fh-L_c,\\
R_{\rm co}&=q+sp/m+Fs^2/(2m)+Z_c-cL_c.
\end{aligned}\tag{59}\] The copy time and shear are predetermined by the
split and do not depend on the unknown force. Both descriptions use the
coordinate and momentum of the same body at the same final time \(h\).
Their distributions differ because their physical interventions differ;
no cross-protocol posterior monotonicity is inferred from (56).

The identities \(\alpha u+\beta h=s\) and
\(\alpha u^2+\beta h^2-s^2=\alpha\beta v^2\) give the exact joint
coupling \[\begin{aligned}
(Q_f,P_f,R_c)
 &=(Q_c,P_c,R_{\rm co})-\frac v m L_{\rm rel}(1,0,\beta)
                         +\delta(0,0,1),\\
\delta&=\frac{F\alpha\beta v^2}{2m},\qquad
\nu=\frac{\kappa^2ruv^3}{hm^2},\qquad w=(1,0,\beta),\\
V_f&=V_c+\nu ww^{\sf T}.
\end{aligned}\tag{60}\] The added noise is independent of the entire
coarse vector. Forgetting the relative record has retained its kick in
the projected law; setting that kick to zero would remove the
\(\nu ww^{\sf T}\) term. The floor (55) holds for each physical protocol
and for the fine block \(R_c\). The law comparison below is joint,
rather than a uniform posterior comparison on every exceptional record.

The coarse covariance \(V_c\) is positive definite. Indeed, the
invertible transformation \(T(Q,P,R)=(Q-hP/m,P,R-Q+\alpha v P/m)\) gives
\[TV_cT^{\sf T}
=\operatorname{diag}(A,B,(rh)^{-1})
  +\kappa^2rh\,aa^{\sf T},\qquad
a=(s/m,-1,-c),\qquad Tw=(1,0,-\alpha).\] Consequently, for
\(x=\nu w^{\sf T}V_c^{-1}w\) and \(y=\delta^2(V_c^{-1})_{33}\),
\[x\le\nu(A^{-1}+\alpha^2rh),\qquad y\le rh\delta^2.\]

Completing the Gaussian square gives the exact affinity between the two
joint laws, with \(d=\delta(0,0,1)\), \[\mathcal A
=\frac{(1+x)^{1/4}}{(1+x/2)^{1/2}}
 \exp\left[-\frac18d^{\sf T}
       (V_c+\nu ww^{\sf T}/2)^{-1}d\right].\tag{61}\] This follows by
whitening \(V_c\), rotating its rank-one perturbation to one coordinate,
and integrating the square root of the two Gaussian densities. The
inverse in the exponent is bounded above by \(V_c^{-1}\). Moreover
\[1-\frac{\sqrt{1+x}}{1+x/2}
=\frac{x^2}{4(1+x/2)(1+x/2+\sqrt{1+x})}\le x^2/8.\] Thus
\(1-\mathcal A^2\le x^2/8+y/4\), and §7's affinity inequality gives
\(\operatorname{TV}\le\frac12\sqrt{x^2/2+y}\). Since
\(\max\alpha\beta^3=27/256\) and \(\max\alpha^3\beta^3=1/64\), define
\[\begin{aligned}
X_h&=\frac{\kappa^2r}{m^2}
       \left[\frac{27h^3}{256A}+\frac{rh^4}{64}\right],\\
Y_h&=\frac{729rF^2h^5}{262144m^2},\\
\left\|\mathcal L(Q_f,P_f,R_c)
       -\mathcal L(Q_c,P_c,R_{\rm co})\right\|_{\rm TV}
 &\le\frac12\sqrt{X_h^2/2+Y_h}.
\end{aligned}\tag{62}\] This is \(O(h^{5/2})\) for bounded \(F\) and
fixed \(A,m,r,\kappa\), and \(O(h^3)\) at \(F=0\), uniformly over the
split. The split maxima follow by differentiating \(\alpha(1-\alpha)^3\)
and by \(\alpha\beta\le1/4\). Here \(r\) has units
\((\text{length}^2\text{time})^{-1}\), \(\nu\) has position-variance
units, \(\delta\) has position units, and \(X_h,Y_h\) are dimensionless.
As either daughter length tends to zero, the residual in (60) vanishes.
At \(\kappa=0\) only the force-dependent mean defect remains.

The comparator's time and shear depend on the split. Equation (62) does
not yet define an associative apparatus rule over many depths. Summing
it over cells would also require a uniform local preparation bound and a
valid conditional comparison after earlier records.

\hypertarget{a-complete-record-information-budget-for-arbitrary-partitions}{%
\subsubsection{8.6 A complete-record information budget for arbitrary
partitions}\label{a-complete-record-information-budget-for-arbitrary-partitions}}

For any fixed, force-independent partition \(0=t_0<\cdots<t_n=\tau\),
\(h_i=t_i-t_{i-1}\), use independent fresh memories with variances
\((1/(rh_i),\kappa^2rh_i)\) and unit copy pulses at the right endpoints.
Keep every physical terminal memory coordinate. The same calculation as
(57) gives \[\begin{aligned}
R_i&=q+t_i p/m+Ft_i^2/(2m)
       -\frac1m\sum_{j<i}(t_i-t_j)\Lambda_j+Z_i,\\
V_R&\succeq\operatorname{diag}((rh_i)^{-1}),\\
\mathcal I_R&\le\frac r{4m^2}\sum_i h_i t_i^4
                      \le\frac{r\tau^5}{4m^2}.
\end{aligned}\tag{63}\] The first formula includes every earlier copy
kick. Each \(Z_i\) is independent of the body and all initial momenta,
so the covariance lower bound follows by adding their diagonal noise
covariance. The record covariance is force-independent and its mean
derivative is \(t_i^2/(2m)\); inverse ordering proves the information
bound. Later copies leave earlier stored coordinates unchanged, as in
§8.4. The terminal vector therefore satisfies (54), regardless of the
number of cells. Its complete force sensitivity obeys, for all \(F,F'\),
\[\|P_F^R-P_{F'}^R\|_{\rm TV}
\le\frac{|F-F'|}{2}\sqrt{\mathcal I_R}
\le\frac{|F-F'|\sqrt r\,\tau^{5/2}}{4m}.\tag{64}\] The first bound is
the Gaussian affinity calculation with equal covariances. It includes
the whole terminal coordinate record and every function of that record,
so relative coordinate channels have not been omitted. Additional
apparatus revealing initial \(Z_i\), reuse of their conjugates, or
adaptive timing require a new calculation.

The total initial precision is fixed by construction:
\(\sum_i rh_i=r\tau\). Hamiltonian mechanics does not select this
preparation schedule or a universal upper bound on \(r\). Within the
stated family, terminal storage and a uniform record-continuity budget
are derived, and the posterior floor follows from the positive joint
preparation scale. Equation (64) remains true at \(\kappa=0\) with the
same imprecision law. Thus continuity in this resource class does not
force positive action. Letting \(r\) grow without bound leaves the
uniform class and is a separate physical-resource question.

The following calculation tests three-cell blocking in both
parenthesizations, carrying relative modes to the terminal projection.
It identifies the timing and noise data the coarse channel retains.

\hypertarget{an-associative-descriptor-with-the-physical-kicks-retained}{%
\subsubsection{8.7 An associative descriptor with the physical kicks
retained}\label{an-associative-descriptor-with-the-physical-kicks-retained}}

Keep §8.6's fixed, force-independent partition and jointly Gaussian
product preparation. The body is centred initially, with \(A,B>0\) and
\(AB\ge\kappa^2\). Allow \(\kappa=0\) in the classical formulas. Define
\(Y=\sum_i h_iR_i\) at the terminal cut. Partial sums of this expression
are bookkeeping for projections of the physical memory coordinates;
every copy pulse is still part of the dynamics. This definition grants
no extra apparatus for passively storing an earlier numerical read.

During one cell of duration \(h\), drift under the deterministic force,
apply its right-endpoint copy, and add that memory coordinate with
weight \(h\). For \(X=(Q,P,Y)^{\sf T}\) the exact update is
\[\begin{aligned}
X'&=M_hX+Ff_h+\zeta_h,\\
M_h&=\begin{pmatrix}1&h/m&0\\0&1&0\\h&h^2/m&1\end{pmatrix},
&f_h&=\begin{pmatrix}h^2/(2m)\\h\\h^3/(2m)\end{pmatrix},\\
\zeta_h&=(0,-\Lambda,hZ)^{\sf T},
&N_h&=\operatorname{Cov}\zeta_h
       =\operatorname{diag}(0,\kappa^2rh,h/r).
\end{aligned}\tag{65}\] The new pair is independent of all earlier
variables. Thus a block's Gaussian channel is described by \((M,f,N)\),
and successive blocks obey \[\begin{aligned}
(M_2,f_2,N_2)\circ(M_1,f_1,N_1)
={}&(M_2M_1,\ M_2f_1+f_2,\ M_2N_1M_2^{\sf T}+N_2),\\
N_{321}={}&M_3M_2N_1M_2^{\sf T}M_3^{\sf T}
             +M_3N_2M_3^{\sf T}+N_3.
\end{aligned}\tag{66}\] Multiplication and the displayed noise sum give
exactly the same \(M,f,N\) in both three-cell parenthesizations. In
particular, each earlier momentum kick propagates through the later
position and record updates. The accumulated record of any subblock has
weight equal to its duration when converted to an average, so its
terminal projection also associates. These descriptors specify the
projected law of the declared physical copy schedule. Their composite
generally differs from the single-cell descriptor \((M_H,f_H,N_H)\) at
the total duration: the composite retains the daughter timing and noise
data. Arbitrary triples are not asserted to have a one-memory
Hamiltonian realization.

\hypertarget{the-three-cell-relative-residue-and-the-necessary-coarse-data}{%
\subsubsection{8.8 The three-cell relative residue and the necessary
coarse
data}\label{the-three-cell-relative-residue-and-the-necessary-coarse-data}}

For a block of duration \(H\), put \(\omega_i=h_i/H\) and define
\[\begin{aligned}
s&=\sum_i\omega_it_i,& k_i&=H-t_i,\\
b_j&=\sum_{i>j}\omega_i(t_i-t_j),&
\Gamma&=\sum_j\omega_jb_j,\\
\sigma_t^2&=\sum_i\omega_it_i^2-s^2,&
Z_c&=\sum_i\omega_iZ_i,\qquad L_c=\sum_i\Lambda_i.
\end{aligned}\] The pair \((Z_c,L_c)\) is canonical, saturated and
Gaussian, with variances \(((rH)^{-1},\kappa^2rH)\). Resolve every
initial momentum as \(\Lambda_i=\omega_iL_c+\Lambda_i^\perp\). Then
\[\operatorname{Cov}(\Lambda_i^\perp,\Lambda_j^\perp)
=\kappa^2rH(\omega_i\delta_{ij}-\omega_i\omega_j).
\tag{67}\] Its covariance with \(L_c\) vanishes, and all memory momenta
are independent of \(Z_c\) and the initial body. Joint Gaussianity
therefore makes the relative vector independent of the entire coarse
vector. It is degenerate in the direction \(\sum_i\Lambda_i^\perp=0\);
no inverse on that null direction is used.

As in §8.5, shear the effective memory to \(Z_c-\Gamma L_c/m\), copy at
\(s\), and drift to \(H\). These controls depend only on the partition.
Let \((Q_c,P_c,R_{\rm co})\) be this physical comparator. Summing (63)
gives the exact coupling \[\begin{aligned}
(Q_f,P_f,Y/H)
 &=(Q_c,P_c,R_{\rm co})-(U/m,0,V/m)
                         +(0,0,F\sigma_t^2/(2m)),\\
U&=\sum_i k_i\Lambda_i^\perp,\qquad
V=\sum_i b_i\Lambda_i^\perp,\\
\operatorname{Cov}(U,V)
 &=\kappa^2rH
 \begin{pmatrix}
 \operatorname{Var}_\omega k&\operatorname{Cov}_\omega(k,b)\\
 \operatorname{Cov}_\omega(k,b)&\operatorname{Var}_\omega b
 \end{pmatrix}.
\end{aligned}\tag{68}\] Here \(\operatorname{Var}_\omega\) denotes the
finite weighted variance with weights \(\omega_i\). The deterministic
difference is the force response of the physical schedule, not a
force-dependent correction which the unknown-force apparatus is assumed
able to apply.

For three equal cells of length \(\ell\), \[\begin{aligned}
k&=(2\ell,\ell,0),&b&=(\ell,\ell/3,0),\\
s&=2\ell,&\Gamma&=4\ell/9,\qquad\sigma_t^2=2\ell^2/3,\\
\operatorname{Cov}(U,V)
 &=\kappa^2r\ell^3
       \begin{pmatrix}2&1\\1&14/27\end{pmatrix},&
\det\operatorname{Cov}(U,V)&=\kappa^4r^2\ell^6/27.
\end{aligned}\tag{69}\] For example
\(\operatorname{Var}_\omega k=2\ell^2/3\),
\(\operatorname{Cov}_\omega(k,b)=\ell^2/3\), and
\(\operatorname{Var}_\omega b=14\ell^2/81\). Multiplication by
\(\kappa^2rH=3\kappa^2r\ell\) yields (69).

At \(\kappa>0\) this independent residue has rank two in the \((Q,Y/H)\)
plane. Even if the deterministic means are matched, the prescribed
saturated effective pair plus one independent scalar entering additively
along a fixed vector supplies a covariance increment of rank at most
one. It cannot reproduce (69). This excludes that specific truncated
rule; correlated re-preparations and other physical coarse apparatus
have not been excluded. A nonlinear vector-valued function of one random
scalar can have covariance rank two and is also outside this exclusion.
Equations (65)--(66) retain the full necessary timing, force-response
and noise data. At \(\kappa=0\) the residue vanishes and the
timing-dependent mean difference remains.

\hypertarget{a-partition-independent-limit-of-the-projected-terminal-channel}{%
\subsubsection{8.9 A partition-independent limit of the projected
terminal
channel}\label{a-partition-independent-limit-of-the-projected-terminal-channel}}

Fix \(\tau,m,r,A,B,\kappa\) and the deterministic force \(F\). The
partition may be arbitrary; set \(\delta=\max_i h_i\). Put
\(s_j=\tau-t_j\) and \[\begin{aligned}
S_1&=\sum_i h_it_i
       =\frac{\tau^2}2+\frac12\sum_i h_i^2,\\
S_2&=\sum_i h_it_i^2,\qquad
0\le S_2-\frac{\tau^3}3\le\tau^2\delta,\\
B_j&=\sum_{i>j}h_i(t_i-t_j)
       =\frac{s_j^2}2+e_j,\qquad
e_j=\frac12\sum_{i>j}h_i^2\le\frac{\delta s_j}2.
\end{aligned}\tag{70}\] For the last identity apply the \(S_1\) identity
to the subpartition after \(t_j\). For \(S_2\), the right-endpoint sum
exceeds the polynomial area \(\tau^3/3\); on each cell its excess is at
most \(\tau h_i^2\), since the slope of \(t^2\) is at most \(2\tau\).
Summing gives the bound. These are finite partition identities and
elementary polynomial estimates.

Repeated composition in (66), starting with \(Y=0\), is consequently
\[\begin{aligned}
M_\pi&=\begin{pmatrix}1&\tau/m&0\\0&1&0\\
                       \tau&S_1/m&1\end{pmatrix},&
f_\pi&=(\tau^2/(2m),\tau,S_2/(2m))^{\sf T},\\
N_\pi&=\kappa^2r\sum_jh_jv_jv_j^{\sf T}
                    +\operatorname{diag}(0,0,\tau/r),&
v_j&=(-s_j/m,-1,-B_j/m)^{\sf T}.
\end{aligned}\tag{71}\] Define a candidate three-dimensional Gaussian
law \(P_*\) by replacing \(S_1,S_2\) in \(M_\pi,f_\pi\) by
\(\tau^2/2,\tau^3/3\), and setting \[N_*=\kappa^2r
\begin{pmatrix}
\tau^3/(3m^2)&\tau^2/(2m)&\tau^4/(8m^2)\\
\tau^2/(2m)&\tau&\tau^3/(6m)\\
\tau^4/(8m^2)&\tau^3/(6m)&\tau^5/(20m^2)
\end{pmatrix}
+\operatorname{diag}(0,0,\tau/r).\tag{72}\] The first matrix is the
polynomial Gram matrix of \((-s/m,-1,-s^2/(2m))\) over
\(0\le s\le\tau\). Thus it is nonnegative. Together with the independent
initial body covariance and the last term, it gives a positive definite
terminal covariance. This defines a finite terminal Gaussian channel
without assuming a limiting mechanical curve or a stochastic process at
all times.

\textbf{Proposition 8 (projected terminal-channel refinement).} Let
\(P_\pi=\mathcal L(Q_f,P_f,Y)\) for the physical schedule. Define
\[\begin{aligned}
a&=\frac{\tau}{m\sqrt A},&b&=\frac1{\sqrt B},&
c&=\sqrt{\frac r\tau}\frac{\tau^2}{2m},\\
L&=\sqrt{a^2+b^2+c^2},&J&=\sqrt{a^2+4c^2},\\
\beta&=\sqrt{\frac{rB}\tau}\frac{\tau\delta}{2m},&
d&=c\delta/\tau,\\
\varepsilon(\delta)
 &=2\beta+\beta^2+
    \kappa^2r\left[LJ\delta+\tau(2Ld+d^2)\right],&
z(\delta)&=|F|c\delta.
\end{aligned}\] Whenever \(\varepsilon(\delta)\le1/2\),
\[\|P_\pi-P_*\|_{\rm TV}
\le\left[\frac34\varepsilon(\delta)^2+
                       \frac13z(\delta)^2\right]^{1/2}.
\tag{73}\] All constants are independent of the number of cells and
their relative lengths. The bound is \(O(\delta)\) at fixed parameters,
uniformly over bounded \(F\). Two partitions are Cauchy by the triangle
inequality, with the sum of their bounds (73). This proves refinement
independence of the body and the single weighted terminal projection in
this apparatus family.

\emph{Proof.} Apply the invertible, dimensionless coordinate change
\[W=\left(\frac{Q-\tau P/m}{\sqrt A},\frac P{\sqrt B},
 \sqrt{\frac r\tau}\left[Y-\tau Q+\frac{\tau^2P}{2m}\right]
 \right)^{\sf T}.\] For the candidate, the initial body and summed
coordinate noise give covariance \(I\), and the memory momenta add a
nonnegative matrix. Thus \(\widehat V_*\succeq I\). For a finite
partition, the former contribution is
\((I+\beta_\pi E_{32})(I+\beta_\pi E_{32})^{\sf T}\), where
\(\beta_\pi=\sqrt{rB/\tau}(S_1-\tau^2/2)/m\le\beta\). Its difference
from \(I\) has norm at most \(2\beta+\beta^2\).

A kick at \(t_j\) transforms to
\[w(t_j)+(0,0,-\sqrt{r/\tau}\,e_j/m)^{\sf T},\qquad
w(t)=\left(\frac t{m\sqrt A},-\frac1{\sqrt B},
                 -\sqrt{\frac r\tau}\frac{t^2}{2m}\right)^{\sf T}.\] On
\([0,\tau]\), \(\|w\|\le L\), \(\|w'\|\le J/\tau\), and the error vector
has norm at most \(d\). The norm of the derivative of \(ww^{\sf T}\) is
at most \(2LJ/\tau\). Its right-endpoint quadrature error on cell \(i\)
is therefore at most \(LJh_i^2/\tau\). Summing gives at most
\(LJ\delta\). The error vector changes each outer product by at most
\(2Ld+d^2\). Equation (71) now proves
\(\|\widehat V_\pi-\widehat V_*\|\le\varepsilon(\delta)\). The only mean
difference in these coordinates is \(F\sqrt{r/\tau}(S_2-\tau^3/3)/(2m)\)
in the third component; (70) bounds its norm by \(z(\delta)\).

Whiten further by \(\widehat V_*^{-1/2}\). Its norm is at most one. The
relative covariance eigenvalues are \(1+x_j\),
\(|x_j|\le\varepsilon\le1/2\), and the mean difference has norm at most
\(z\). The squared Gaussian affinity is \[\mathcal A^2=\prod_{j=1}^3
 \frac{\sqrt{1+x_j}}{1+x_j/2}
 \exp\left[-\frac14\mu^{\sf T}
                         (I+E/2)^{-1}\mu\right].\] For \(|x|\le1/2\),
rationalizing the square root gives \[1-\frac{\sqrt{1+x}}{1+x/2}
=\frac{x^2}{4(1+x/2)(1+x/2+\sqrt{1+x})}\le x^2/4.\] The denominator is
at least \(17/4>4\), using \(1+x/2\ge3/4\) and \(\sqrt{1+x}\ge2/3\).
Also \(I+E/2\succeq(3/4)I\), so the exponential's deficit is at most
\(z^2/3\). The deficit of a product of factors in \([0,1]\) is bounded
by the sum of their deficits. Hence
\(1-\mathcal A^2\le3\varepsilon^2/4+z^2/3\). The affinity inequality of
§7 gives (73). \(\square\)

Here \([r]=(\text{length}^2\text{time})^{-1}\), \([\kappa]\) is action,
and \(Y\) has length-times-time units. The constants \(a,b,c,L,J,d\)
have inverse-momentum units; \(\beta,\varepsilon,z\) are dimensionless.
Every term in (71)--(72) has the covariance units of its coordinate
pair. The transformed variables and the statistical bound are
dimensionless.

For \(\kappa>0\), each finite \(P_\pi\) satisfies terminal posterior
closure (55), since \(Y\) is a linear projection of commuting physical
terminal coordinates. Its record variance is at least \(\tau/r\); the
conditional covariance is therefore continuous at \(P_*\), and the limit
retains \(\det\Sigma_{(Q,P)\mid Y}\ge\kappa^2\). At \(\kappa=0\),
(65)--(73) still hold with all momentum kicks zero; convergence and
continuity in this fixed-precision family select no positive action
scale.

Section 8.10 tests the retained storage geometry under a physically
realized adaptive-gain controller, including its conjugate recoil.
Equation (73) controls one projected channel; extending it to the full
growing record requires compatible record spaces and conditional
comparison estimates.

\textbf{Referee verdict (GPT-6 Astra, 2026-09-30): ACCEPT at this finite
Gaussian level.} The descriptor, kick signs, rank-two covariance and
constants in (73) were re-derived. The exclusion in §8.8 is restricted
to an independent additive scalar along a fixed vector, and the mean
bound retains the centred initial-body assumption. The posterior limit
uses covariance convergence and a positive record variance, rather than
TV convergence alone.

\hypertarget{a-physical-adaptive-gain-controller-with-complete-record-closure}{%
\subsubsection{8.10 A physical adaptive-gain controller with
complete-record
closure}\label{a-physical-adaptive-gain-controller-with-complete-record-closure}}

Retain §8.6's force-independent partition and product Gaussian
preparation, including the centred initial body. For each pulse choose a
force-independent \(C^2\) function \(g_i(R_1,\ldots,R_{i-1})\), with
\(|g_i|\le G<\infty\) and bounded first derivatives for each finite
protocol. The gains are dimensionless. At \(t_i\) use the integrated
Hamiltonian generator \(\mathcal G_i=g_i(R_{<i})Q_i\Lambda_i\). The
arguments are the actual stored coordinates of earlier memories, which
commute with the current body. This is a supplied nonlinear mechanical
coupling, with its full canonical recoil included.

During this pulse \(Q_i\), \(R_{<i}\) and the fresh momentum
\(\Lambda_i\) stay fixed. Its exact flow is \[\begin{aligned}
P&\longmapsto P-g_i\Lambda_i,\qquad
Z_i\longmapsto R_i=Z_i+g_iQ_i,\\
\Pi_j&\longmapsto\Pi_j-(\partial_jg_i)Q_i\Lambda_i
                         \quad(j<i).
\end{aligned}\tag{74}\] Here \(\Pi_j\) denotes the current conjugate of
stored coordinate \(R_j\). The symbol \(\Lambda_j\) in all subsequent
terminal formulas denotes its initial, fresh pre-copy momentum, rather
than its later recoiled conjugate. Old memories are otherwise frozen;
their conjugates remain in the full mechanical state, and are isolated
from later body couplings and readouts. Later generators use their
coordinates directly. Thus the controller uses physical storage
throughout. The \(C^2\) regularity makes the recoil-inclusive flow a
differentiable canonical map; bounded first derivatives and gains keep
its finite-protocol moments finite. No implementation resource bound for
such couplings is derived here.

Summing the body impulses gives exactly \[\begin{aligned}
R_i&=Z_i+g_i(R_{<i})\left[q+t_ip/m+Ft_i^2/(2m)
 -\frac1m\sum_{j<i}(t_i-t_j)g_j(R_{<j})\Lambda_j\right],\\
Q_\tau&=q+\tau p/m+F\tau^2/(2m)
 -\frac1m\sum_j(\tau-t_j)g_j(R_{<j})\Lambda_j,\\
P_\tau&=p+F\tau-\sum_jg_j(R_{<j})\Lambda_j.
\end{aligned}\tag{75}\] Every delivered kick survives. The complete
declared coordinate record is \(R=(R_1,\ldots,R_n)\); old conjugate
readouts and terminal-body readouts would require additional measuring
couplings and a new law.

\textbf{Proposition 9 (physical adaptive-gain closure).} Under these
apparatus assumptions, conditionally on each complete terminal value
\(R=a\), the body law is Gaussian and satisfies
\[\Sigma_{(Q_\tau,P_\tau)\mid R=a}
             +i\kappa\Omega_{\rm body}\succeq0.\tag{76}\] In particular
its determinant is at least \(\kappa^2\) at \(\kappa>0\). Let \(P_F^R\)
be the adaptive coordinate-record law. Then \[\begin{aligned}
J(a)&\le C=\frac{rG^2\tau^5}{4m^2},\\
\mathcal A(F,F')&=E_{(F+F')/2}
                    e^{-(F-F')^2J(R)/8},\\
\|P_F^R-P_{F'}^R\|_{\rm TV}
&\le\sqrt{1-e^{-(F-F')^2C/4}}
 \le\frac{|F-F'|G\sqrt r\,\tau^{5/2}}{4m}.
\end{aligned}\tag{77}\] The cap is uniform in the finite partition and
all controllers in the bounded-gain class. The continuity calculation
remains valid at \(\kappa=0\).

\emph{Proof.} At fixed \(a\), freeze the numerical gains at
\(g_i=g_i(a_{<i})\). Write \(C_g,L_g\) for the ordinary linear record
and terminal-body matrices of that predetermined-gain copy protocol, and
put \(c_{g,i}=g_it_i^2/(2m)\),
\(V_g=\operatorname{Cov}(C_gZ_{\rm init})\). Each record has its own
independent \(Z_i\) with coefficient one, so
\(V_g\succeq D=\operatorname{diag}((rh_i)^{-1})\succ0\).

To check the adaptive density rather than assume Gaussianity, fix the
latent variables \(U=(q,p,\Lambda_1,\ldots,\Lambda_n)\). The map
\((Z_i)\mapsto(R_i)\) in (75) is triangular with diagonal derivatives
one. Its inverse is \(Z_i=a_i-g_i(a_{<i})Q_i(U;g_{<i}(a),F)\). If
\(\varphi_i\) is the initial density of \(Z_i\), integrate against the
Gaussian latent law \(\mu_U\): \[p_F(a)=\int\prod_i
 \varphi_i\bigl(a_i-g_i(a_{<i})Q_i(U;g_{<i}(a),F)\bigr)
 \,d\mu_U(U)
=\varphi_{Fc_g,V_g}(a)\big|_{g=g(a)}.\tag{78}\] This is the frozen-gain
Gaussian density evaluated at the adaptive path. Its derivation
establishes normalization and includes all gain dependence; no
gain-derivative Jacobian is present. The full adaptive law is generally
non-Gaussian.

The same integral before integrating \(U\) identifies its conditional
law at \(a\) as precisely the frozen protocol's Gaussian posterior. At
that fixed record the terminal body in (75) is the same affine function
\(L_gZ_{\rm init}\), with its force mean. For the auxiliary frozen
linear protocol, copies and drifts give (54) for \(L_g,C_g\).
Proposition 7 therefore proves (76). This applies its covariance theorem
to the frozen protocol; the full nonlinear terminal state need not
satisfy a Gaussian covariance closure claim.

At the same path \(a\), \(g(a)\) and \(V_g\) are identical under both
force hypotheses. Completing the Gaussian square in (78) gives
\[\sqrt{p_F(a)p_{F'}(a)}
=p_{(F+F')/2}(a)
           e^{-(F-F')^2J(a)/8},\qquad
J(a)=c_g^{\sf T}V_g^{-1}c_g.\] Inverse ordering and the gain bound imply
\[J(a)\le c_g^{\sf T}D^{-1}c_g
=\frac r{4m^2}\sum_i h_ig_i(a_{<i})^2t_i^4
\le C.\] Integration proves the affinity identity and lower bound
\(\mathcal A\ge e^{-(F-F')^2C/8}\). The affinity inequality and
\(1-e^{-x}\le x\) give (77). \(\square\)

Frozen-protocol Gaussian innovations along the observed path are also
the actual adaptive conditional kernels along that path. Their
information sums to \(J(a)\), consistent with (41)--(44). Conditional
kernels compose and their branchwise information budgets add. When
composing blocks, retain the body's conditional mean, covariance and
inherited force sensitivity: applying a fresh local-duration
\(\tau_{\rm block}^5\) cap would drop force information already in the
incoming body. Equation (77) is the valid global-duration bound. It also
controls every measurable function of the declared full record.

\textbf{Referee verdict (GPT-6 Astra, 2026-09-30): ACCEPT with the
\(C^2\) and momentum-notation qualifications stated above.} The pulse
recoil, triangular Jacobian, posterior substitution, affinity identity
and constant were re-derived. The zero branch and the conditional
composition scope are retained.

Fixed \(r,G\) are supplied resource parameters; (77) selects no positive
action. Adaptive timing and compatibility of different gain rules under
refinement remain open. The next Newton lemma is an inserted-cell
comparison for a prescribed smooth policy depending on the physical
weighted past record, with both terminal-body dynamics and controller
recoil retained. Dependencies are (74)--(78), a common physical coarse
projection, and a conditional noise/moment bound. Abandon a rule if its
feedback uses passive numerical storage or if its projected-law error
loses a mesh-uniform bound. The initial positive preparation law,
universality and radiation calibration remain separate obligations.

\hypertarget{consequence-for-state}{%
\subsection{9. Consequence for STATE}\label{consequence-for-state}}

The Newton task advances from testing stationary radiation to deriving
an admissible recording law. Shared Gaussian driving, including its
cross correlations and radiation damping, allows the two-pointer
counterexample at every fixed finite cutoff when classical retained
readouts and the stated controls are available. The next physical
premise to justify is a restriction such as (20)--(22) for reusable
pointers and their memories. The referee's analysis narrows it: the
background's prior already obeys (22), so the premise is a bound on
every terminal classical readout, a restriction on observation
equivalent at \(\kappa=\hbar/2\) to Gaussian quantum measurement theory;
an electrodynamic instrument constraint would only relocate it. The
conditional scale identification and the unregulated momentum problem
retain the scope stated in the revised SED link. The September 29
checkpoints sharpen the recording law to conditional innovation
(32)--(35), extend complete-record continuity and composition to
adaptive Gaussian kernels (41)--(44), and calculate the innovation
replacement cost (46). The sister formalization's norm obstacle yields a
finite harmonic refinement bound; (49) states its extra statistical
budget. Section 8 supplies a Hamiltonian copy/monitor realization and
terminal closure from the initial joint phase-covariance restriction.
Sections 8.4--8.10 give physical copy refinement, the complete-record
information budget (63), associative projected-channel composition (66),
and the partition-independent terminal Gaussian law at rate (73). The
three-cell residue (69) determines why the coarse channel retains more
noise data than the prescribed one-pair/scalar rule. A physical
adaptive-gain controller retains its conjugate recoil and gives
posterior closure (76) and a global complete-record continuity bound
(77). Next Newton: the inserted-cell comparison of a prescribed physical
feedback policy. Alternate next to the gauge normalized interaction
forest. Positivity, universality and radiation calibration remain
separate obligations.
\end{document}
